% =====================================================================
% extensions.tex
% ---------------------------------------------------------------------
% Additive extensions to conf.tex addressing reviewer concerns:
%
%   (a) Piecewise-stationary lower bound: extends Theorem~\ref{thm:spsc_lb}
%       from K=1 to K segments via a per-segment application of the same
%       single-segment construction. Yields the matching
%       Omega(c^{1/3} K^{1/3} T^{2/3}) rate.
%
%   (b) Sharpened d-dependence in the probe term: replaces the
%       worst-case envelope analysis (giving C_sub = Theta(d)) with a
%       variance-adaptive matrix Bernstein argument (giving
%       C_sub = Theta(sqrt d)), tightening Theorem~\ref{thm:spsc_regret}'s
%       probe term from O~(d^{2/3} K^{1/3} T^{2/3}) to
%       O~(d^{1/3} K^{1/3} T^{2/3}).
%
% Both are PURELY ADDITIVE: existing Theorems~\ref{thm:spsc_regret},
% \ref{thm:spsc_adaptive}, \ref{thm:spsc_lb} are unchanged. The
% experimental section is unaffected.
%
% INTEGRATION:
%   * Add `% =====================================================================
% extensions.tex
% ---------------------------------------------------------------------
% Additive extensions to conf.tex addressing reviewer concerns:
%
%   (a) Piecewise-stationary lower bound: extends Theorem~\ref{thm:spsc_lb}
%       from K=1 to K segments via a per-segment application of the same
%       single-segment construction. Yields the matching
%       Omega(c^{1/3} K^{1/3} T^{2/3}) rate.
%
%   (b) Sharpened d-dependence in the probe term: replaces the
%       worst-case envelope analysis (giving C_sub = Theta(d)) with a
%       variance-adaptive matrix Bernstein argument (giving
%       C_sub = Theta(sqrt d)), tightening Theorem~\ref{thm:spsc_regret}'s
%       probe term from O~(d^{2/3} K^{1/3} T^{2/3}) to
%       O~(d^{1/3} K^{1/3} T^{2/3}).
%
% Both are PURELY ADDITIVE: existing Theorems~\ref{thm:spsc_regret},
% \ref{thm:spsc_adaptive}, \ref{thm:spsc_lb} are unchanged. The
% experimental section is unaffected.
%
% INTEGRATION:
%   * Add `% =====================================================================
% extensions.tex
% ---------------------------------------------------------------------
% Additive extensions to conf.tex addressing reviewer concerns:
%
%   (a) Piecewise-stationary lower bound: extends Theorem~\ref{thm:spsc_lb}
%       from K=1 to K segments via a per-segment application of the same
%       single-segment construction. Yields the matching
%       Omega(c^{1/3} K^{1/3} T^{2/3}) rate.
%
%   (b) Sharpened d-dependence in the probe term: replaces the
%       worst-case envelope analysis (giving C_sub = Theta(d)) with a
%       variance-adaptive matrix Bernstein argument (giving
%       C_sub = Theta(sqrt d)), tightening Theorem~\ref{thm:spsc_regret}'s
%       probe term from O~(d^{2/3} K^{1/3} T^{2/3}) to
%       O~(d^{1/3} K^{1/3} T^{2/3}).
%
% Both are PURELY ADDITIVE: existing Theorems~\ref{thm:spsc_regret},
% \ref{thm:spsc_adaptive}, \ref{thm:spsc_lb} are unchanged. The
% experimental section is unaffected.
%
% INTEGRATION:
%   * Add `% =====================================================================
% extensions.tex
% ---------------------------------------------------------------------
% Additive extensions to conf.tex addressing reviewer concerns:
%
%   (a) Piecewise-stationary lower bound: extends Theorem~\ref{thm:spsc_lb}
%       from K=1 to K segments via a per-segment application of the same
%       single-segment construction. Yields the matching
%       Omega(c^{1/3} K^{1/3} T^{2/3}) rate.
%
%   (b) Sharpened d-dependence in the probe term: replaces the
%       worst-case envelope analysis (giving C_sub = Theta(d)) with a
%       variance-adaptive matrix Bernstein argument (giving
%       C_sub = Theta(sqrt d)), tightening Theorem~\ref{thm:spsc_regret}'s
%       probe term from O~(d^{2/3} K^{1/3} T^{2/3}) to
%       O~(d^{1/3} K^{1/3} T^{2/3}).
%
% Both are PURELY ADDITIVE: existing Theorems~\ref{thm:spsc_regret},
% \ref{thm:spsc_adaptive}, \ref{thm:spsc_lb} are unchanged. The
% experimental section is unaffected.
%
% INTEGRATION:
%   * Add `\input{extensions}` to conf.tex just before `\end{document}`
%     (i.e., after Section J / Table tab:published_comparison).
%   * Optionally add an entry to the appendix TOC (line ~847 of conf.tex):
%       \tocentry{K}{app:extensions}{Extensions: piecewise LB and sharpened d-dependence}
% =====================================================================

\section{Extensions: piecewise lower bound and sharpened \texorpdfstring{$d$}{d}-dependence}
\label{app:extensions}

This section supplements (without modifying) Theorems~\ref{thm:spsc_regret}--\ref{thm:spsc_lb}
with two additive results: a piecewise-stationary lower bound that
matches the upper-bound rate in $K$ (\S\ref{app:ext_piecewise_lb}),
and a sharpened $d$-dependence in the probe term via variance-adaptive
matrix Bernstein (\S\ref{app:ext_sharpened_d}). The original theorems
remain in force; the results below are stated as \emph{independent}
sharpenings.

\subsection{Piecewise-stationary lower bound}
\label{app:ext_piecewise_lb}

Theorem~\ref{thm:spsc_lb} establishes the LB only for $K=1$ (single
segment of length $T$). The $K$-segment extension follows by applying
the single-segment bound conditionally inside each segment, exploiting
the fact that the per-segment hidden direction $v_k$ in the
construction is independent of all previous segments. No new technical
tools beyond Theorem~\ref{thm:spsc_lb} itself are required.

\begin{theorem}[Piecewise lower bound; extension of Theorem~\ref{thm:spsc_lb} to $K\ge 1$]
\label{thm:spsc_lb_piecewise}
Let $\Pi$ be the policy class of Theorem~\ref{thm:spsc_lb} (bounded
exploit actions $\|x_t\|\le R_\cA$, isotropic Gaussian probes
$u_t\sim\cN(0,I_d)$). Fix $K\ge 1$ and segment lengths
$\{\ell_k\}_{k=1}^K$ with $\sum_k\ell_k=T$. There exist constants
$c_0,d_0,\ell_0>0$ depending only on $R_\cA,\sigma_\varepsilon,c$
(in fact, the same constants as Theorem~\ref{thm:spsc_lb}) such that
whenever $\min_k\ell_k\ge\ell_0$ and $d\ge d_0\,(\min_k\ell_k)^{1/3}$,
\begin{equation}
\inf_{\pi\in\Pi}\sup_\nu\ \E^\pi_\nu[\DynReg_T^{(c)}]
\;\ge\; c_0\,c^{1/3}(R_\cA\sigma_\varepsilon)^{2/3}\sum_{k=1}^K\ell_k^{2/3}.
\label{eq:lb_piecewise}
\end{equation}
For uniform segments $\ell_k=T/K$ this yields
\begin{equation}
\inf_{\pi\in\Pi}\sup_\nu\ \E^\pi_\nu[\DynReg_T^{(c)}]
\;\ge\; c_0\,c^{1/3}(R_\cA\sigma_\varepsilon)^{2/3}\,K^{1/3}T^{2/3},
\label{eq:lb_piecewise_uniform}
\end{equation}
matching the probe-acquisition rate of Theorem~\ref{thm:spsc_regret}
in both $K$ and $T$ (with the same residual $d$-prefactor gap discussed
in Remark~\ref{rem:dim_dep}, sharpened in
\S\ref{app:ext_sharpened_d} below).
\end{theorem}

\begin{proof}
\emph{Construction.} Let $v_1,\dots,v_K$ be i.i.d.\ uniform on
$\mathbb S^{d-1}$, drawn independently of the policy. For
per-segment scales $\alpha_k>0$ to be specified, define the instance
$\nu(\vec v):\theta_t = \alpha_k v_k$ for $t\in\cI_k$. Let
$\nu_0:\theta_t\equiv 0$ denote the global null. The per-segment
``zero-on-segment-$k$'' alternative is
$\nu^{(-k)}(\vec v):\theta_t = \alpha_{k'}v_{k'}\mathbf 1\{k'\ne k\}$
for $t\in\cI_{k'}$ — i.e., identical to $\nu(\vec v)$ except on
segment $k$ where $\theta_t=0$.

\emph{Per-segment KL.} The trajectory laws under $\nu(\vec v)$ and
$\nu^{(-k)}(\vec v)$ agree on segments $k'\ne k$ and differ only on
segment $k$. The sequential KL chain rule therefore gives
\[
\mathrm{KL}\bigl(\mathbb P_{\nu^{(-k)}(\vec v)}\,\bigl\|\,\mathbb P_{\nu(\vec v)}\bigr)
\;=\; \frac{\alpha_k^2}{2\sigma_\varepsilon^2}\,
\E^\pi_{\nu^{(-k)}(\vec v)}\!\Bigl[\sum_{t\in\cI_k}(x_t^\top v_k)^2\Bigr].
\]
Under $\nu^{(-k)}(\vec v)$, on segment $k$ the parameter is $0$, so
the segment-$k$ marginal trajectory is independent of $v_k$.
Consequently $v_k$ is independent of the actions $\{x_t\}_{t\in\cI_k}$
under $\nu^{(-k)}(\vec v)$, and the trace-identity argument of
Lemma~\ref{lem:trace_kl_conf} applies verbatim with horizon $\ell_k$,
probe count $m_k$, and exploit count $n_k=\ell_k-m_k$:
\[
\E_{v_k}\E^\pi_{\nu^{(-k)}(\vec v)}\!\Bigl[\sum_{t\in\cI_k}(x_t^\top v_k)^2\Bigr]
\;\le\; m_k + R_\cA^2 n_k/d \;=:\;Q_k,
\]
where the average is over $v_k\sim\mathrm{Unif}(\mathbb S^{d-1})$. As in
Lemma~\ref{lem:trace_kl_conf}, the existence of an extremal direction
gives a deterministic $v_k^\star\in\mathbb S^{d-1}$ for which the
inequality holds with $v_k^\star$ replacing the average. Setting
$\alpha_k^2 := \sigma_\varepsilon^2/(4 Q_k)$ yields
\[
\mathrm{KL}\bigl(\mathbb P_{\nu^{(-k)}(\vec v^\star)}\,\bigl\|\,\mathbb P_{\nu(\vec v^\star)}\bigr)
\;\le\;\frac{\alpha_k^2 Q_k}{2\sigma_\varepsilon^2}\;=\;\frac{1}{8},
\]
hence Pinsker gives per-segment TV $\le 1/4$.

\emph{Per-segment regret transfer.} The optimal action under
$\nu(\vec v^\star)$ on segment $k$ is $x_t^\star=R_\cA v_k^\star$ with
reward $\alpha_k R_\cA$; under $\nu^{(-k)}(\vec v^\star)$ the
parameter on segment $k$ is $0$ so any action is optimal, and the
segment-$k$ regret of the policy is identically $0$. The
arm-regret-vs.-alignment manipulation of \S\ref{app:proof_lb}
(equation~\eqref{eq:reg_to_alignment_conf}), applied to the segment-$k$
exploit rounds, gives
\[
\E^\pi_{\nu(\vec v^\star)}\!\Bigl[\sum_{t\in\cI_k\cap\cT_{\mathrm{ex}}}\Delta_t\Bigr]
\;\ge\;\alpha_k n_k R_\cA\;-\;\alpha_k\,\E^\pi_{\nu(\vec v^\star)}[f_k],\quad
f_k:=\sum_{t\in\cI_k\cap\cT_{\mathrm{ex}}}|v_k^{\star\top}x_t|.
\]
TV transfer from $\nu(\vec v^\star)$ to $\nu^{(-k)}(\vec v^\star)$
with $\|f_k\|_\infty\le n_k R_\cA$ and segment-$k$
$d_{\mathrm{TV}}\le 1/4$:
\[
\E^\pi_{\nu(\vec v^\star)}[f_k]\;\le\;\E^\pi_{\nu^{(-k)}(\vec v^\star)}[f_k]+\tfrac12\,n_k R_\cA.
\]
Under $\nu^{(-k)}(\vec v^\star)$, by Cauchy--Schwarz and the
trace-identity bound applied on segment $k$,
$\E^\pi_{\nu^{(-k)}(\vec v^\star)}[f_k]\le R_\cA n_k/\sqrt d$.
Combining (with $d\ge 16$ so $1/\sqrt d\le 1/4$):
\begin{equation}
\E^\pi_{\nu(\vec v^\star)}\!\Bigl[\sum_{t\in\cI_k\cap\cT_{\mathrm{ex}}}\Delta_t\Bigr]
\;\ge\;\tfrac14\,\alpha_k n_k R_\cA
\;=\;\frac{R_\cA\sigma_\varepsilon\,n_k}{8\sqrt{Q_k}}.
\label{eq:per_seg_regret_lb}
\end{equation}

\emph{Per-segment costed regret.} Adding the segment-$k$ probe cost
$cm_k$,
\[
\E^\pi_{\nu(\vec v^\star)}[\DynReg^{(c)}_k]
\;\ge\; cm_k + \frac{R_\cA\sigma_\varepsilon\,n_k}{8\sqrt{Q_k}},
\]
with the same functional form as the single-segment bound
\eqref{eq:dynreg_lower_conf}. The Regime-I/Regime-II optimization of
\S\ref{app:proof_lb} applies verbatim with $T\to\ell_k$, yielding
\[
\E^\pi_{\nu(\vec v^\star)}[\DynReg^{(c)}_k]
\;\ge\;c_0\,c^{1/3}(R_\cA\sigma_\varepsilon)^{2/3}\ell_k^{2/3}
\]
in the regime $\ell_k\ge\ell_0$ and $d\ge d_0\,\ell_k^{1/3}$, with
$c_0,d_0,\ell_0$ as in Theorem~\ref{thm:spsc_lb}.

\emph{Summing across segments.} Total dynamic regret is the sum of
per-segment costed regrets (probes and exploit rounds partition by
segment):
\[
\E^\pi_{\nu(\vec v^\star)}[\DynReg^{(c)}_T]
\;=\;\sum_{k=1}^K\E^\pi_{\nu(\vec v^\star)}[\DynReg^{(c)}_k]
\;\ge\;c_0\,c^{1/3}(R_\cA\sigma_\varepsilon)^{2/3}\sum_{k=1}^K\ell_k^{2/3}.
\]
Taking $\sup_\nu\ge$ this specific instance yields
\eqref{eq:lb_piecewise}; specialising to $\ell_k=T/K$ gives
$\sum_k\ell_k^{2/3}=K\,(T/K)^{2/3}=K^{1/3}T^{2/3}$, hence
\eqref{eq:lb_piecewise_uniform}.\qed
\end{proof}

\begin{remark}[The argument uses only Theorem~\ref{thm:spsc_lb}'s tools]
\label{rem:lb_piecewise_tools}
The proof reuses the trace-identity Lemma~\ref{lem:trace_kl_conf} and
the Le~Cam transfer of \S\ref{app:proof_lb} per segment. The only new
ingredient is the observation that the segment-$k$ KL between
$\nu(\vec v^\star)$ and $\nu^{(-k)}(\vec v^\star)$ is exactly the
single-segment KL with parameters $(\ell_k, m_k, n_k)$, because the
two laws agree outside segment $k$. No multi-hypothesis Fano, no
sliding-window reduction, and no new concentration is invoked.
\end{remark}

\begin{remark}[Adaptive probe schedules]
\label{rem:lb_piecewise_adaptive}
The bound \eqref{eq:per_seg_regret_lb} is in terms of the policy's
\emph{realised} segment-$k$ probe count $m_k$ and exploit count $n_k$,
which may be adaptive (chosen based on history through $\tau_{k-1}$).
The per-segment Regime-I/Regime-II optimization minimizes
$cm_k+R_\cA\sigma_\varepsilon n_k/(8\sqrt{Q_k})$ over $m_k\in[0,\ell_k]$;
the policy's adaptive choice of $m_k$ can only weakly increase this
quantity. Hence the LB applies under any (possibly adaptive)
per-segment probe allocation.
\end{remark}

\begin{remark}[Combined LB/UB picture]
\label{rem:combined_lb_ub}
Theorem~\ref{thm:spsc_lb_piecewise} combined with
Theorem~\ref{thm:spsc_regret} yields the matching pair
\[
\Omega(c^{1/3}\,K^{1/3}T^{2/3})
\;\le\;\text{(probe-acquisition regret)}\;\le\;
\tilO(d^{2/3}\,K^{1/3}T^{2/3}),
\]
tight in $K$ and $T$. The remaining gap is the $d$-prefactor on the
upper-bound side; \S\ref{app:ext_sharpened_d} narrows that gap from
$d^{2/3}$ to $d^{1/3}$ via a sharper concentration argument, leaving
a residual $d^{1/3}$ open question.
\end{remark}

\subsection{Sharpened \texorpdfstring{$d$}{d}-dependence via variance-adaptive Bernstein}
\label{app:ext_sharpened_d}

Theorem~\ref{thm:matrix_bernstein_conf} (matrix Freedman) bounds
$\|\widehat M_k-\bar M_k^{\mathrm{probe}}-\widetilde B\|_\op$ using the
worst-case envelope $\|\tilde X_t\|_\op\le 2R_X$ with
$R_X=\Theta(L^2)=\Theta(d\log T)$. The leading term
$2R_X\sqrt{\log(2d/\delta)/m_k}$ then yields
$C_{\mathrm{sub}}=8R_X/\lambda_{\min}=\Theta(d\log T/\lambda_{\min})$
in Corollary~\ref{cor:projector_conf}, and the constant
$B = 4 C_{\mathrm{sub}} S_w R_\cA\sqrt{\log(2d/\delta)} = \Theta(d)$
propagates through the proof of Theorem~\ref{thm:spsc_regret} to give
the $\tilO(d^{2/3}K^{1/3}T^{2/3})$ probe-acquisition rate.

The envelope is loose. For isotropic Gaussian $u_t$,
$\|u_t\|^2$ concentrates near $d$ with $\sqrt d$ fluctuations, so the
\emph{predictable second moment} of $\tilde G_t$ in operator norm is
$\Theta(d)$ per round, not $\Theta(d^2)$ as the envelope suggests.
The variance-adaptive matrix Bernstein
inequality~\citep[Theorem~7.7.1]{tropp2015introduction} exploits this
gap and yields a sharper concentration. The algorithmic structure of
SPSC (probe-then-exploit interleaving, projector recovery, projected
ridge-UCB) is unchanged; only the recommended probe rate $m_k$ is
updated by a $d^{1/3}$ factor
(from $m_k^\star\asymp d^{2/3}\ell_k^{2/3}$ to
$m_k^\star\asymp d^{1/3}\ell_k^{2/3}$; see
Theorem~\ref{thm:spsc_regret_sharp} and Remark~\ref{rem:no_algo_change}).

\begin{lemma}[Predictable second moment of the lifted probe sample]
\label{lem:lifted_variance}
Assume $d\ge 2$, $\|\theta_t\|\le S_w$ a.s., and that the noise has a
bounded fourth moment $\E[\varepsilon_t^4\mid\cH_{t-1}]\le 3\sigma_\varepsilon^4$
(satisfied by Gaussian $\varepsilon_t\sim\cN(0,\sigma_\varepsilon^2)$).
Let $u_t\sim\cN(0,I_d)$ on probe rounds (truncation by an indicator
$\mathbf 1\{\|u_t\|\le L\}$ only decreases the second moment in the
PSD order, so the bound below applies; if probes are instead
distributed $u_t\mid\{\|u_t\|\le L\}$, the same bound holds up to a
$1+o(1)$ factor as $L\to\infty$). Define $G_t := \mathcal K^{-1}(s_t u_t u_t^\top)$
and $\tilde X_t := G_t\mathbf 1\{\cA_t\}-\E[G_t\mathbf 1\{\cA_t\}\mid\cH_{t-1}]$.
Then there exists $\sigma_*^2$ depending only on $S_w,\sigma_\varepsilon,\hat\sigma$
(and not on $d$) such that
\begin{equation}
\bigl\|\E[\tilde X_t^2\mid\cH_{t-1}]\bigr\|_\op
\;\le\;\sigma_*^2\,d.
\label{eq:lifted_var}
\end{equation}
Explicitly, $\sigma_*^2 \le C\,(S_w^4 + \sigma_\varepsilon^4 + \hat\sigma^4)$
with $C=100$ (using the exact trace-expansion bound established in the
proof; the value of $C$ does not affect the leading-order rate).
\end{lemma}
\begin{proof}
Centering only reduces the second moment in the PSD order:
$\E[\tilde X_t^2\mid\cH_{t-1}]\preceq\E[G_t^2\mathbf 1\{\cA_t\}\mid\cH_{t-1}]\preceq\E[G_t^2\mid\cH_{t-1}]$.
We bound $\|\E[G_t^2\mid\cH_{t-1}]\|_\op$ by conditioning further on
$\theta_t$ and proving a bound uniform in $\theta_t$ with
$\|\theta_t\|\le S_w$. By tower-property the same bound then applies
to $\E[\,\cdot\mid\cH_{t-1}]$.

By Lemma~\ref{lem:K_inverse},
\(
G_t = \tfrac{s_t}{2}\,u_t u_t^\top - \tfrac{s_t\|u_t\|^2}{2(d{+}2)}\,I.
\)
Squaring and using $(uu^\top)^2 = \|u\|^2 uu^\top$:
\begin{equation}
G_t^2 = \frac{d}{4(d+2)}\,s_t^2\|u_t\|^2\, u_t u_t^\top
       \;+\;\frac{s_t^2\,\|u_t\|^4}{4(d+2)^2}\,I.
\label{eq:Gt_squared}
\end{equation}

\emph{Leading-in-$d$ block.}
Conditional on $(\cH_{t-1},\theta_t)$, define $e_1 := \theta_t/\|\theta_t\|$
(any unit vector if $\theta_t=0$). This $e_1$ is now deterministic, so
the decomposition $u_t = z_t e_1 + u_{t,\perp}$ is well-defined with
$z_t := u_t^\top e_1\sim\cN(0,1)$ and $u_{t,\perp}$ independent of
$z_t$, distributed as $\cN(0,I-e_1e_1^\top)$ on the $(d{-}1)$-dim
subspace $e_1^\perp$. Crucially, $s_t = (z_t\|\theta_t\|+\varepsilon_t)^2-\hat\sigma^2$
depends on $u_t$ only through $z_t$, so $s_t$ is independent of
$u_{t,\perp}$ given $(\cH_{t-1},\theta_t)$.

Expand $u_tu_t^\top = z_t^2 e_1e_1^\top + z_t(e_1 u_{t,\perp}^\top + u_{t,\perp}e_1^\top) + u_{t,\perp}u_{t,\perp}^\top$
and $\|u_t\|^2 = z_t^2 + \|u_{t,\perp}\|^2$. The cross terms $z_t e_1 u_{t,\perp}^\top$
vanish under $\E[\cdot\mid\cH_{t-1},\theta_t]$ by symmetry of
$u_{t,\perp}$. The surviving contributions split into:
\begin{align*}
\E\bigl[s_t^2\|u_t\|^2\,z_t^2\,e_1e_1^\top \,\big|\,\cH_{t-1},\theta_t\bigr]
&= \Bigl(\E[s_t^2 z_t^4] + (d{-}1)\,\E[s_t^2 z_t^2]\Bigr)\,e_1e_1^\top, \\
\E\bigl[s_t^2\|u_t\|^2\,u_{t,\perp}u_{t,\perp}^\top \,\big|\,\cH_{t-1},\theta_t\bigr]
&= \E[s_t^2 z_t^2]\,(I-e_1e_1^\top) + \E[s_t^2]\,(d{+}1)(I-e_1e_1^\top),
\end{align*}
where conditional expectations are abbreviated, and we used:
\begin{itemize}[leftmargin=1.5em,nosep]
\item $\E[u_{t,\perp}u_{t,\perp}^\top]=I-e_1e_1^\top$ (covariance);
\item $\E[\|u_{t,\perp}\|^2]=d-1$ ($(d{-}1)$-dim chi-square);
\item $\E[\|u_{t,\perp}\|^2\,u_{t,\perp}u_{t,\perp}^\top]=(d+1)(I-e_1e_1^\top)$
  (Isserlis on the $(d{-}1)$-dim subspace, valid for $d\ge 2$);
\item independence of $u_{t,\perp}$ from $s_t$ given $(\cH_{t-1},\theta_t)$
  splits the joint expectation in the second display.
\end{itemize}

\emph{Bounding the conditional moments of $s_t$.} Write
$\xi_t := z_t\|\theta_t\|+\varepsilon_t$, so $s_t = \xi_t^2-\hat\sigma^2$ and
$s_t^2 = \xi_t^4 - 2\hat\sigma^2\xi_t^2 + \hat\sigma^4$. Given
$\theta_t$, $\xi_t$ has variance $\|\theta_t\|^2+\sigma_\varepsilon^2\le S_w^2+\sigma_\varepsilon^2$.
Using the assumed fourth-moment bound $\E[\varepsilon_t^4\mid\cH_{t-1}]\le 3\sigma_\varepsilon^4$
and Gaussian moments $\E[z_t^k]=k!!$ for $k\le 8$ (in particular
$\E[z_t^8]=105$ — needed because $s_t^2 z_t^4$ contains a $z_t^8$
contribution from $\xi_t^4 z_t^4 = (z_t\|\theta_t\|+\varepsilon_t)^4 z_t^4$),
straightforward expansion gives
\[
\E[s_t^2\mid\cH_{t-1},\theta_t]\le C_0\,D,\quad
\E[s_t^2 z_t^2\mid\cH_{t-1},\theta_t]\le C_0\,D,\quad
\E[s_t^2 z_t^4\mid\cH_{t-1},\theta_t]\le C_0\,D,
\]
with $D := S_w^4+\sigma_\varepsilon^4+\hat\sigma^4$ and a universal
constant $C_0$ (the dominant contribution is
$\E[\xi_t^4 z_t^4]\le 105\,S_w^4 + 90\,S_w^2\sigma_\varepsilon^2+9\sigma_\varepsilon^4
\le 320\,D$ via $\E[z_t^8]=105$ and AM-GM, so $C_0\le 320$ suffices).

\emph{Combining the blocks.} The parallel-block coefficient is
$A_1 := \E[s_t^2 z_t^4]+(d-1)\E[s_t^2 z_t^2]\le d\,C_0\,D$, and the
perpendicular-block coefficient is
$A_2 := \E[s_t^2 z_t^2]+(d+1)\E[s_t^2]\le (d+2)\,C_0\,D$. Hence
$\|\E[s_t^2\|u_t\|^2 u_tu_t^\top\mid\cH_{t-1},\theta_t]\|_\op
\le \max(A_1,A_2)\le (d+2)\,C_0\,D$. Multiplying by the
$G_t^2$-prefactor $d/(4(d+2))$ from \eqref{eq:Gt_squared} yields a
leading-block contribution
$\le (d/(4(d+2)))\cdot (d+2)\,C_0\,D = (C_0/4)\,d\,D$ exactly.

\emph{Trace correction (exact bound).} The second term of
\eqref{eq:Gt_squared} contributes
$\E[s_t^2\|u_t\|^4]/(4(d+2)^2)\cdot I$. Expanding
$\|u_t\|^4 = z_t^4 + 2z_t^2\|u_{t,\perp}\|^2 + \|u_{t,\perp}\|^4$ and
using $\E[\|u_{t,\perp}\|^4]=(d-1)(d+1)$,
\[
\E[s_t^2\|u_t\|^4\mid\cH_{t-1},\theta_t]
\;\le\;C_0\,D\,\bigl(1+2(d-1)+(d^2-1)\bigr)
\;=\;C_0\,D\,(d^2+2d-2).
\]
Divided by $4(d+2)^2$, the trace contribution to
$\|\E[G_t^2\mid\cH_{t-1},\theta_t]\|_\op$ is at most
$C_0\,D\,(d^2+2d-2)/(4(d+2)^2)$.

\emph{Conclusion (exact prefactor accounting).} Combining the
leading-block contribution $\le (C_0/4)\,d\,D$ with the trace
correction yields, after dividing by $d\,D$,
\[
\frac{\|\E[G_t^2\mid\cH_{t-1},\theta_t]\|_\op}{d\,D}
\;\le\;\frac{C_0}{4}
\;+\;\frac{C_0\,(d^2+2d-2)}{4d\,(d+2)^2}.
\]
With $C_0=320$, the second term is monotone decreasing in $d$ for
$d\ge 2$: at $d=2$, $C_0\cdot 6/(8\cdot 16)=15$; at $d=4$,
$C_0\cdot 22/(16\cdot 36)\approx 12.22$; as $d\to\infty$, the term
approaches $0$. Combined with the constant first term $C_0/4=80$,
the per-$dD$ bound is $\le 80+15=95$ at $d=2$, $\le 92.22$ at $d=4$,
and $\to 80$ as $d\to\infty$ — uniformly $\le 100$. Hence
$\sigma_*^2 = C\,D$ with $C = 100$ suffices for all $d\ge 2$.
Taking conditional expectation over $\theta_t$ preserves the bound
by linearity and operator-norm convexity, yielding
\eqref{eq:lifted_var}.
\end{proof}

\begin{theorem}[Variance-adaptive concentration; sharpening of Theorem~\ref{thm:matrix_bernstein_conf}]
\label{thm:matrix_bernstein_var}
Under the assumptions of Theorem~\ref{thm:matrix_bernstein_conf}
together with the additional hypotheses of Lemma~\ref{lem:lifted_variance}
($d\ge 2$ and $\E[\varepsilon_t^4\mid\cH_{t-1}]\le 3\sigma_\varepsilon^4$;
truncation by indicator $\mathbf 1\{\|u_t\|\le L\}$), with
probability at least $1-2\delta$ (jointly over the truncation event of
Lemma~\ref{lem:G_bound_conf} and the Bernstein event),
\begin{equation}
\bigl\|\widehat M_k-\bar M_k^{\mathrm{probe}}-\widetilde B\bigr\|_\op
\;\le\;\sigma_*\sqrt{\frac{2\,d\,\log(2d/\delta)}{m_k}}
\;+\;\frac{4 R_X\log(2d/\delta)}{3\,m_k}
\;+\;\|\Theta_k\|_\op,
\label{eq:bern_var}
\end{equation}
where $R_X = \Theta(L^2)=\Theta(d\log T)$ is the worst-case envelope of
Lemma~\ref{lem:G_bound_conf} and $\sigma_*^2$ is the variance constant
of Lemma~\ref{lem:lifted_variance}.
\end{theorem}
\begin{proof}
This is matrix Bernstein \citep[Theorem~7.7.1]{tropp2015introduction}
applied to the centered MDS $\{\tilde X_t\}_{t\in\cT_k}$, with:
\begin{itemize}[leftmargin=1.5em,nosep]
\item a.s.\ envelope $L_X := \|\tilde X_t\|_\op\le 2R_X$ (Lemma~\ref{lem:G_bound_conf}, plus centering);
\item predictable variance
$\bigl\|\sum_{t\in\cT_k}\E[\tilde X_t^2\mid\cH_{t-1}]\bigr\|_\op\le m_k\,\sigma_*^2\,d$
(Lemma~\ref{lem:lifted_variance}, summed).
\end{itemize}
Tropp's matrix Bernstein with these inputs gives, with probability
$\ge 1-\delta$,
\[
\Bigl\|\sum_{t\in\cT_k}\tilde X_t\Bigr\|_\op
\;\le\;\sqrt{2 m_k\sigma_*^2 d\log(2d/\delta)} + \frac{2L_X\log(2d/\delta)}{3}
\;\le\;\sqrt{2 m_k\sigma_*^2 d\log(2d/\delta)} + \frac{4 R_X\log(2d/\delta)}{3},
\]
using $L_X\le 2R_X$ in the linear-envelope term. Dividing by $m_k$
and combining with the truncation-bias offset $\|\Theta_k\|_\op$ via
the same transfer step as in Theorem~\ref{thm:matrix_bernstein_conf}
yields \eqref{eq:bern_var}. The two constituent events (truncation
and Bernstein) each have mass $\ge 1-\delta$; union bound gives
$1-2\delta$.
\end{proof}

\begin{corollary}[Sharpened projector concentration; sharpening of Corollary~\ref{cor:projector_conf}]
\label{cor:projector_var}
Under the assumptions of Corollary~\ref{cor:projector_conf} together
with the additional hypotheses of Lemma~\ref{lem:lifted_variance}
($d\ge 2$, $\E[\varepsilon_t^4\mid\cH_{t-1}]\le 3\sigma_\varepsilon^4$,
indicator-truncated probes), with probability at least $1-2\delta$,
\begin{equation}
\varepsilon_k = \|\widehat P_k - P_k^\star\|_\op
\;\le\; C_{\mathrm{sub}}^\star\,\sqrt{\frac{d\,\log(2d/\delta)}{m_k}}
\;+\;\frac{16\,R_X\,\log(2d/\delta)}{3\,\lambda_{\min}\,m_k}
\;+\;\Delta_\sigma\;+\;\frac{4\|\Theta_k\|_\op}{\lambda_{\min}},
\quad
C_{\mathrm{sub}}^\star \;:=\;\frac{4\sqrt{2}\,\sigma_*}{\lambda_{\min}}.
\label{eq:proj_bound_var}
\end{equation}
The linear-in-$1/m_k$ term is dominated by the leading
$\sqrt{d/m_k}$ term whenever $m_k\gtrsim R_X^2\log(2d/\delta)/(\sigma_*^2 d)$,
which evaluates to $m_k\gtrsim d\log^2 T\cdot\log(2d/\delta) = \tilde\Theta(d)$
(with $R_X=\Theta(d\log T)$ from Lemma~\ref{lem:G_bound_conf}). The
$\tilde\Theta$ hides polylog factors of $T$ and $d/\delta$.
The sharpened constant $C_{\mathrm{sub}}^\star = \Theta(1/\lambda_{\min})$
(in the variance scale) gives leading-term scaling
$O(\sqrt{d/m_k})$, replacing the original
$C_{\mathrm{sub}}=\Theta(d)/\lambda_{\min}$ leading scaling
$O(d/\sqrt{m_k})$ — a factor-of-$\sqrt d$ improvement in the
$m_k$-dependence at any fixed $d$.
\end{corollary}
\begin{proof}
Davis--Kahan applied to \eqref{eq:bern_var}, with the leading
$\sigma_*\sqrt{2d\log(2d/\delta)/m_k}$ replacing the original
$2R_X\sqrt{\log(2d/\delta)/m_k}=\Theta(d\log T\sqrt{\log(2d/\delta)/m_k})$.
The other terms (linear $1/m_k$ Bernstein term, truncation bias,
variance-misspecification floor $\Delta_\sigma$) are unchanged.
\end{proof}

\begin{theorem}[Sharpened probe term; sharpening of Theorem~\ref{thm:spsc_regret}]
\label{thm:spsc_regret_sharp}
In addition to the assumptions of Theorem~\ref{thm:spsc_regret},
assume the conditions of Lemma~\ref{lem:lifted_variance}: $d\ge 2$
and $\E[\varepsilon_t^4\mid\cH_{t-1}]\le 3\sigma_\varepsilon^4$
(satisfied by Gaussian noise; truncation is by indicator
$\mathbf 1\{\|u_t\|\le L\}$, so $\mathcal K^{-1}$ is the full-Gaussian
inverse used throughout). Run Alg.~\ref{alg:spsc} with the sharpened
probe allocation $m_k = \lceil(B^\star\ell_k/(2A))^{2/3}\rceil$, where
\[
A := 2R_\cA S_w + c, \qquad
B^\star := 4 C_{\mathrm{sub}}^\star\,S_w\,R_\cA\,\sqrt{d\,\log(2d/\delta)}
\;=\;\Theta(\sqrt d).
\]
Assume the per-segment regime
$\ell_k\ge \ell^\star = \tilde\Theta(d)$ for all $k\in[K]$
(explicitly, $\ell^\star = C_\ell\,d\,\mathrm{polylog}(T,d/\delta)$
for a universal constant $C_\ell$, where the polylog absorbs the
$\log(2d/\delta)$ and $\log T$ factors of Corollary~\ref{cor:projector_var};
when segments are balanced, this is equivalent to
$T\ge K\ell^\star = \tilde\Theta(Kd)$).
Then, with probability $\ge 1-\delta$,
\begin{equation}
\DynReg_T^{(c)}
\;=\;\tilO\!\bigl(r\sqrt{KT}\bigr)
\;+\;\tilO\!\bigl(d^{1/3}\,K^{1/3}\,T^{2/3}\bigr)
\;+\;O(WV)\;+\;O(T\Delta_\sigma).
\label{eq:main_bound_sharp}
\end{equation}
The probe-term prefactor improves from $d^{2/3}$ to $d^{1/3}$.
\end{theorem}
\begin{proof}
The proof of Theorem~\ref{thm:spsc_regret} (\S\ref{app:proof_main})
uses $\varepsilon_k\le C_{\mathrm{sub}}\sqrt{\log(2d/\delta)/m_k}+\Delta_\sigma$
inside the step (i.c) tradeoff
\[
2R_\cA S_w\,\varepsilon_k\,n_k \;\le\; B\,n_k/\sqrt{m_k} + 2R_\cA S_w\Delta_\sigma\,n_k,
\quad B = 4 C_{\mathrm{sub}}\,S_w\,R_\cA\sqrt{\log(2d/\delta)}.
\]
Corollary~\ref{cor:projector_var} replaces this with
\[
\varepsilon_k \le C_{\mathrm{sub}}^\star\sqrt{\frac{d\log(2d/\delta)}{m_k}}
+\frac{16 R_X\log(2d/\delta)}{3\lambda_{\min}m_k}+\Delta_\sigma.
\]
Multiplying by $2R_\cA S_w n_k$, the leading variance term becomes
$B^\star n_k/\sqrt{m_k}$ with $B^\star = \Theta(\sqrt d)$ as displayed,
and the linear-Bernstein term contributes
$32 R_\cA S_w R_X \log(2d/\delta)\,n_k/(3\lambda_{\min}m_k) = \tilO(d\,n_k/m_k)$.

\emph{Threshold derivation.} From Corollary~\ref{cor:projector_var},
the linear-Bernstein contribution to the regret tradeoff is
$\tilO(d\,n_k/m_k)$ and the leading variance contribution is
$\tilO(\sqrt d\,n_k/\sqrt{m_k})$ (both up to polylog factors of
$T$ and $d/\delta$, which differ between the two terms but get
absorbed into $\tilde\Theta$ uniformly below). Their ratio at the sharpened optimum
$m_k^\star\asymp d^{1/3}\ell_k^{2/3}$ is
\[
\frac{\text{leading}}{\text{linear}}
\;\asymp\;
\frac{\sqrt{d/m_k^\star}}{d/m_k^\star}
\;=\;\sqrt{m_k^\star/d}
\;\asymp\;(\ell_k/d)^{1/3}.
\]
For domination $\ge 1$ (modulo polylog), $\ell_k\gtrsim d$ up to
polylog factors. We absorb these into $\ell^\star = \tilde\Theta(d)$;
the explicit form $C_\ell\,d\,\mathrm{polylog}(T,d/\delta)$ handles
both the $R_X = O(d\log T)$ envelope and the $\log(2d/\delta)$
Bernstein-tail factor.

\emph{Per-segment optimization.} In the regime $\ell_k\ge\ell^\star$,
the linear term is dominated, and
$\min_{m_k}\bigl(A m_k+B^\star\ell_k/\sqrt{m_k}\bigr)$
is achieved at $m_k^\star = (B^\star\ell_k/(2A))^{2/3}$ with value
\[
\Theta\bigl(A^{1/3}(B^\star)^{2/3}\bigr)\,\ell_k^{2/3} = \Theta(d^{1/3})\,\ell_k^{2/3}
\]
per segment. The leading-to-linear ratio at $m_k^\star$ is
$(B^\star\ell_k/\sqrt{m_k^\star})/(d\,n_k/m_k^\star)
\asymp (\ell_k/d)^{1/3}\to\infty$ under $\ell_k\ge\ell^\star$, confirming
domination.

\emph{Summing over segments.} Under the theorem's hypothesis
$\ell_k\ge\ell^\star$ for all $k\in[K]$, the per-segment
contributions sum via Jensen
$\sum_k\ell_k^{2/3}\le K^{1/3}T^{2/3}$ to
$\tilO(d^{1/3}K^{1/3}T^{2/3})$. (For balanced segments
$\ell_k=T/K$, the hypothesis reduces to $T\ge K\ell^\star=\tilde\Theta(Kd)$,
satisfied for $T$ large at any fixed $K,d$.)

All other steps of the proof of Theorem~\ref{thm:spsc_regret} (the
windowed-ridge in-subspace argument, the drift summation, the
variance-misspecification floor) carry over verbatim because they do
not involve $C_{\mathrm{sub}}$.\qed
\end{proof}

\begin{remark}[Algorithm structure unchanged; recommended probe rate updated]
\label{rem:no_algo_change}
Theorem~\ref{thm:spsc_regret_sharp} preserves the algorithmic
structure of Alg.~\ref{alg:spsc} (probe-then-exploit interleaving,
projector recovery via $\widehat M_k$, windowed projected ridge-UCB);
only the recommended probe rate is updated:
$m_k^\star\asymp d^{1/3}\ell_k^{2/3}$ replaces the original
$m_k^\star\asymp d^{2/3}\ell_k^{2/3}$, a $d^{1/3}$-factor reduction.
The empirical results of \S\ref{sec:experiments} were obtained with
the original schedule, so they correspond to a (theoretically)
suboptimal — but consistent — probe rate; reported regret is therefore
an upper bound on what the sharpened schedule would achieve.
The empirical $d$-scaling slope $\approx-0.16$ in Figure~\ref{fig:d-scaling}
is closer to the sharpened $d^{1/3}=+0.33$ prediction than to the
original $d^{2/3}=+0.67$, providing post-hoc empirical support for
the variance-adaptive analysis.
\end{remark}

\begin{remark}[Combined LB/UB picture, after both extensions]
\label{rem:combined_after_ext}
Theorem~\ref{thm:spsc_lb_piecewise} and
Theorem~\ref{thm:spsc_regret_sharp} together yield
\[
\Omega\!\bigl(c^{1/3}\,K^{1/3}T^{2/3}\bigr)
\;\le\;\text{(probe-acquisition regret of SPSC)}\;\le\;
\tilO\!\bigl(d^{1/3}\,K^{1/3}T^{2/3}\bigr),
\]
tight in $K$ and $T$, with a residual $d^{1/3}$ gap that the
empirical $d$-scaling fit (Figure~\ref{fig:d-scaling}) suggests is
also loose; closing it likely requires the subspace-restricted
concentration sketched in Remark~\ref{rem:dim_dep}, which is beyond
the scope of this purely-additive sharpening.
\end{remark}

% =====================================================================
% End extensions.tex
% =====================================================================
` to conf.tex just before `\end{document}`
%     (i.e., after Section J / Table tab:published_comparison).
%   * Optionally add an entry to the appendix TOC (line ~847 of conf.tex):
%       \tocentry{K}{app:extensions}{Extensions: piecewise LB and sharpened d-dependence}
% =====================================================================

\section{Extensions: piecewise lower bound and sharpened \texorpdfstring{$d$}{d}-dependence}
\label{app:extensions}

This section supplements (without modifying) Theorems~\ref{thm:spsc_regret}--\ref{thm:spsc_lb}
with two additive results: a piecewise-stationary lower bound that
matches the upper-bound rate in $K$ (\S\ref{app:ext_piecewise_lb}),
and a sharpened $d$-dependence in the probe term via variance-adaptive
matrix Bernstein (\S\ref{app:ext_sharpened_d}). The original theorems
remain in force; the results below are stated as \emph{independent}
sharpenings.

\subsection{Piecewise-stationary lower bound}
\label{app:ext_piecewise_lb}

Theorem~\ref{thm:spsc_lb} establishes the LB only for $K=1$ (single
segment of length $T$). The $K$-segment extension follows by applying
the single-segment bound conditionally inside each segment, exploiting
the fact that the per-segment hidden direction $v_k$ in the
construction is independent of all previous segments. No new technical
tools beyond Theorem~\ref{thm:spsc_lb} itself are required.

\begin{theorem}[Piecewise lower bound; extension of Theorem~\ref{thm:spsc_lb} to $K\ge 1$]
\label{thm:spsc_lb_piecewise}
Let $\Pi$ be the policy class of Theorem~\ref{thm:spsc_lb} (bounded
exploit actions $\|x_t\|\le R_\cA$, isotropic Gaussian probes
$u_t\sim\cN(0,I_d)$). Fix $K\ge 1$ and segment lengths
$\{\ell_k\}_{k=1}^K$ with $\sum_k\ell_k=T$. There exist constants
$c_0,d_0,\ell_0>0$ depending only on $R_\cA,\sigma_\varepsilon,c$
(in fact, the same constants as Theorem~\ref{thm:spsc_lb}) such that
whenever $\min_k\ell_k\ge\ell_0$ and $d\ge d_0\,(\min_k\ell_k)^{1/3}$,
\begin{equation}
\inf_{\pi\in\Pi}\sup_\nu\ \E^\pi_\nu[\DynReg_T^{(c)}]
\;\ge\; c_0\,c^{1/3}(R_\cA\sigma_\varepsilon)^{2/3}\sum_{k=1}^K\ell_k^{2/3}.
\label{eq:lb_piecewise}
\end{equation}
For uniform segments $\ell_k=T/K$ this yields
\begin{equation}
\inf_{\pi\in\Pi}\sup_\nu\ \E^\pi_\nu[\DynReg_T^{(c)}]
\;\ge\; c_0\,c^{1/3}(R_\cA\sigma_\varepsilon)^{2/3}\,K^{1/3}T^{2/3},
\label{eq:lb_piecewise_uniform}
\end{equation}
matching the probe-acquisition rate of Theorem~\ref{thm:spsc_regret}
in both $K$ and $T$ (with the same residual $d$-prefactor gap discussed
in Remark~\ref{rem:dim_dep}, sharpened in
\S\ref{app:ext_sharpened_d} below).
\end{theorem}

\begin{proof}
\emph{Construction.} Let $v_1,\dots,v_K$ be i.i.d.\ uniform on
$\mathbb S^{d-1}$, drawn independently of the policy. For
per-segment scales $\alpha_k>0$ to be specified, define the instance
$\nu(\vec v):\theta_t = \alpha_k v_k$ for $t\in\cI_k$. Let
$\nu_0:\theta_t\equiv 0$ denote the global null. The per-segment
``zero-on-segment-$k$'' alternative is
$\nu^{(-k)}(\vec v):\theta_t = \alpha_{k'}v_{k'}\mathbf 1\{k'\ne k\}$
for $t\in\cI_{k'}$ — i.e., identical to $\nu(\vec v)$ except on
segment $k$ where $\theta_t=0$.

\emph{Per-segment KL.} The trajectory laws under $\nu(\vec v)$ and
$\nu^{(-k)}(\vec v)$ agree on segments $k'\ne k$ and differ only on
segment $k$. The sequential KL chain rule therefore gives
\[
\mathrm{KL}\bigl(\mathbb P_{\nu^{(-k)}(\vec v)}\,\bigl\|\,\mathbb P_{\nu(\vec v)}\bigr)
\;=\; \frac{\alpha_k^2}{2\sigma_\varepsilon^2}\,
\E^\pi_{\nu^{(-k)}(\vec v)}\!\Bigl[\sum_{t\in\cI_k}(x_t^\top v_k)^2\Bigr].
\]
Under $\nu^{(-k)}(\vec v)$, on segment $k$ the parameter is $0$, so
the segment-$k$ marginal trajectory is independent of $v_k$.
Consequently $v_k$ is independent of the actions $\{x_t\}_{t\in\cI_k}$
under $\nu^{(-k)}(\vec v)$, and the trace-identity argument of
Lemma~\ref{lem:trace_kl_conf} applies verbatim with horizon $\ell_k$,
probe count $m_k$, and exploit count $n_k=\ell_k-m_k$:
\[
\E_{v_k}\E^\pi_{\nu^{(-k)}(\vec v)}\!\Bigl[\sum_{t\in\cI_k}(x_t^\top v_k)^2\Bigr]
\;\le\; m_k + R_\cA^2 n_k/d \;=:\;Q_k,
\]
where the average is over $v_k\sim\mathrm{Unif}(\mathbb S^{d-1})$. As in
Lemma~\ref{lem:trace_kl_conf}, the existence of an extremal direction
gives a deterministic $v_k^\star\in\mathbb S^{d-1}$ for which the
inequality holds with $v_k^\star$ replacing the average. Setting
$\alpha_k^2 := \sigma_\varepsilon^2/(4 Q_k)$ yields
\[
\mathrm{KL}\bigl(\mathbb P_{\nu^{(-k)}(\vec v^\star)}\,\bigl\|\,\mathbb P_{\nu(\vec v^\star)}\bigr)
\;\le\;\frac{\alpha_k^2 Q_k}{2\sigma_\varepsilon^2}\;=\;\frac{1}{8},
\]
hence Pinsker gives per-segment TV $\le 1/4$.

\emph{Per-segment regret transfer.} The optimal action under
$\nu(\vec v^\star)$ on segment $k$ is $x_t^\star=R_\cA v_k^\star$ with
reward $\alpha_k R_\cA$; under $\nu^{(-k)}(\vec v^\star)$ the
parameter on segment $k$ is $0$ so any action is optimal, and the
segment-$k$ regret of the policy is identically $0$. The
arm-regret-vs.-alignment manipulation of \S\ref{app:proof_lb}
(equation~\eqref{eq:reg_to_alignment_conf}), applied to the segment-$k$
exploit rounds, gives
\[
\E^\pi_{\nu(\vec v^\star)}\!\Bigl[\sum_{t\in\cI_k\cap\cT_{\mathrm{ex}}}\Delta_t\Bigr]
\;\ge\;\alpha_k n_k R_\cA\;-\;\alpha_k\,\E^\pi_{\nu(\vec v^\star)}[f_k],\quad
f_k:=\sum_{t\in\cI_k\cap\cT_{\mathrm{ex}}}|v_k^{\star\top}x_t|.
\]
TV transfer from $\nu(\vec v^\star)$ to $\nu^{(-k)}(\vec v^\star)$
with $\|f_k\|_\infty\le n_k R_\cA$ and segment-$k$
$d_{\mathrm{TV}}\le 1/4$:
\[
\E^\pi_{\nu(\vec v^\star)}[f_k]\;\le\;\E^\pi_{\nu^{(-k)}(\vec v^\star)}[f_k]+\tfrac12\,n_k R_\cA.
\]
Under $\nu^{(-k)}(\vec v^\star)$, by Cauchy--Schwarz and the
trace-identity bound applied on segment $k$,
$\E^\pi_{\nu^{(-k)}(\vec v^\star)}[f_k]\le R_\cA n_k/\sqrt d$.
Combining (with $d\ge 16$ so $1/\sqrt d\le 1/4$):
\begin{equation}
\E^\pi_{\nu(\vec v^\star)}\!\Bigl[\sum_{t\in\cI_k\cap\cT_{\mathrm{ex}}}\Delta_t\Bigr]
\;\ge\;\tfrac14\,\alpha_k n_k R_\cA
\;=\;\frac{R_\cA\sigma_\varepsilon\,n_k}{8\sqrt{Q_k}}.
\label{eq:per_seg_regret_lb}
\end{equation}

\emph{Per-segment costed regret.} Adding the segment-$k$ probe cost
$cm_k$,
\[
\E^\pi_{\nu(\vec v^\star)}[\DynReg^{(c)}_k]
\;\ge\; cm_k + \frac{R_\cA\sigma_\varepsilon\,n_k}{8\sqrt{Q_k}},
\]
with the same functional form as the single-segment bound
\eqref{eq:dynreg_lower_conf}. The Regime-I/Regime-II optimization of
\S\ref{app:proof_lb} applies verbatim with $T\to\ell_k$, yielding
\[
\E^\pi_{\nu(\vec v^\star)}[\DynReg^{(c)}_k]
\;\ge\;c_0\,c^{1/3}(R_\cA\sigma_\varepsilon)^{2/3}\ell_k^{2/3}
\]
in the regime $\ell_k\ge\ell_0$ and $d\ge d_0\,\ell_k^{1/3}$, with
$c_0,d_0,\ell_0$ as in Theorem~\ref{thm:spsc_lb}.

\emph{Summing across segments.} Total dynamic regret is the sum of
per-segment costed regrets (probes and exploit rounds partition by
segment):
\[
\E^\pi_{\nu(\vec v^\star)}[\DynReg^{(c)}_T]
\;=\;\sum_{k=1}^K\E^\pi_{\nu(\vec v^\star)}[\DynReg^{(c)}_k]
\;\ge\;c_0\,c^{1/3}(R_\cA\sigma_\varepsilon)^{2/3}\sum_{k=1}^K\ell_k^{2/3}.
\]
Taking $\sup_\nu\ge$ this specific instance yields
\eqref{eq:lb_piecewise}; specialising to $\ell_k=T/K$ gives
$\sum_k\ell_k^{2/3}=K\,(T/K)^{2/3}=K^{1/3}T^{2/3}$, hence
\eqref{eq:lb_piecewise_uniform}.\qed
\end{proof}

\begin{remark}[The argument uses only Theorem~\ref{thm:spsc_lb}'s tools]
\label{rem:lb_piecewise_tools}
The proof reuses the trace-identity Lemma~\ref{lem:trace_kl_conf} and
the Le~Cam transfer of \S\ref{app:proof_lb} per segment. The only new
ingredient is the observation that the segment-$k$ KL between
$\nu(\vec v^\star)$ and $\nu^{(-k)}(\vec v^\star)$ is exactly the
single-segment KL with parameters $(\ell_k, m_k, n_k)$, because the
two laws agree outside segment $k$. No multi-hypothesis Fano, no
sliding-window reduction, and no new concentration is invoked.
\end{remark}

\begin{remark}[Adaptive probe schedules]
\label{rem:lb_piecewise_adaptive}
The bound \eqref{eq:per_seg_regret_lb} is in terms of the policy's
\emph{realised} segment-$k$ probe count $m_k$ and exploit count $n_k$,
which may be adaptive (chosen based on history through $\tau_{k-1}$).
The per-segment Regime-I/Regime-II optimization minimizes
$cm_k+R_\cA\sigma_\varepsilon n_k/(8\sqrt{Q_k})$ over $m_k\in[0,\ell_k]$;
the policy's adaptive choice of $m_k$ can only weakly increase this
quantity. Hence the LB applies under any (possibly adaptive)
per-segment probe allocation.
\end{remark}

\begin{remark}[Combined LB/UB picture]
\label{rem:combined_lb_ub}
Theorem~\ref{thm:spsc_lb_piecewise} combined with
Theorem~\ref{thm:spsc_regret} yields the matching pair
\[
\Omega(c^{1/3}\,K^{1/3}T^{2/3})
\;\le\;\text{(probe-acquisition regret)}\;\le\;
\tilO(d^{2/3}\,K^{1/3}T^{2/3}),
\]
tight in $K$ and $T$. The remaining gap is the $d$-prefactor on the
upper-bound side; \S\ref{app:ext_sharpened_d} narrows that gap from
$d^{2/3}$ to $d^{1/3}$ via a sharper concentration argument, leaving
a residual $d^{1/3}$ open question.
\end{remark}

\subsection{Sharpened \texorpdfstring{$d$}{d}-dependence via variance-adaptive Bernstein}
\label{app:ext_sharpened_d}

Theorem~\ref{thm:matrix_bernstein_conf} (matrix Freedman) bounds
$\|\widehat M_k-\bar M_k^{\mathrm{probe}}-\widetilde B\|_\op$ using the
worst-case envelope $\|\tilde X_t\|_\op\le 2R_X$ with
$R_X=\Theta(L^2)=\Theta(d\log T)$. The leading term
$2R_X\sqrt{\log(2d/\delta)/m_k}$ then yields
$C_{\mathrm{sub}}=8R_X/\lambda_{\min}=\Theta(d\log T/\lambda_{\min})$
in Corollary~\ref{cor:projector_conf}, and the constant
$B = 4 C_{\mathrm{sub}} S_w R_\cA\sqrt{\log(2d/\delta)} = \Theta(d)$
propagates through the proof of Theorem~\ref{thm:spsc_regret} to give
the $\tilO(d^{2/3}K^{1/3}T^{2/3})$ probe-acquisition rate.

The envelope is loose. For isotropic Gaussian $u_t$,
$\|u_t\|^2$ concentrates near $d$ with $\sqrt d$ fluctuations, so the
\emph{predictable second moment} of $\tilde G_t$ in operator norm is
$\Theta(d)$ per round, not $\Theta(d^2)$ as the envelope suggests.
The variance-adaptive matrix Bernstein
inequality~\citep[Theorem~7.7.1]{tropp2015introduction} exploits this
gap and yields a sharper concentration. The algorithmic structure of
SPSC (probe-then-exploit interleaving, projector recovery, projected
ridge-UCB) is unchanged; only the recommended probe rate $m_k$ is
updated by a $d^{1/3}$ factor
(from $m_k^\star\asymp d^{2/3}\ell_k^{2/3}$ to
$m_k^\star\asymp d^{1/3}\ell_k^{2/3}$; see
Theorem~\ref{thm:spsc_regret_sharp} and Remark~\ref{rem:no_algo_change}).

\begin{lemma}[Predictable second moment of the lifted probe sample]
\label{lem:lifted_variance}
Assume $d\ge 2$, $\|\theta_t\|\le S_w$ a.s., and that the noise has a
bounded fourth moment $\E[\varepsilon_t^4\mid\cH_{t-1}]\le 3\sigma_\varepsilon^4$
(satisfied by Gaussian $\varepsilon_t\sim\cN(0,\sigma_\varepsilon^2)$).
Let $u_t\sim\cN(0,I_d)$ on probe rounds (truncation by an indicator
$\mathbf 1\{\|u_t\|\le L\}$ only decreases the second moment in the
PSD order, so the bound below applies; if probes are instead
distributed $u_t\mid\{\|u_t\|\le L\}$, the same bound holds up to a
$1+o(1)$ factor as $L\to\infty$). Define $G_t := \mathcal K^{-1}(s_t u_t u_t^\top)$
and $\tilde X_t := G_t\mathbf 1\{\cA_t\}-\E[G_t\mathbf 1\{\cA_t\}\mid\cH_{t-1}]$.
Then there exists $\sigma_*^2$ depending only on $S_w,\sigma_\varepsilon,\hat\sigma$
(and not on $d$) such that
\begin{equation}
\bigl\|\E[\tilde X_t^2\mid\cH_{t-1}]\bigr\|_\op
\;\le\;\sigma_*^2\,d.
\label{eq:lifted_var}
\end{equation}
Explicitly, $\sigma_*^2 \le C\,(S_w^4 + \sigma_\varepsilon^4 + \hat\sigma^4)$
with $C=100$ (using the exact trace-expansion bound established in the
proof; the value of $C$ does not affect the leading-order rate).
\end{lemma}
\begin{proof}
Centering only reduces the second moment in the PSD order:
$\E[\tilde X_t^2\mid\cH_{t-1}]\preceq\E[G_t^2\mathbf 1\{\cA_t\}\mid\cH_{t-1}]\preceq\E[G_t^2\mid\cH_{t-1}]$.
We bound $\|\E[G_t^2\mid\cH_{t-1}]\|_\op$ by conditioning further on
$\theta_t$ and proving a bound uniform in $\theta_t$ with
$\|\theta_t\|\le S_w$. By tower-property the same bound then applies
to $\E[\,\cdot\mid\cH_{t-1}]$.

By Lemma~\ref{lem:K_inverse},
\(
G_t = \tfrac{s_t}{2}\,u_t u_t^\top - \tfrac{s_t\|u_t\|^2}{2(d{+}2)}\,I.
\)
Squaring and using $(uu^\top)^2 = \|u\|^2 uu^\top$:
\begin{equation}
G_t^2 = \frac{d}{4(d+2)}\,s_t^2\|u_t\|^2\, u_t u_t^\top
       \;+\;\frac{s_t^2\,\|u_t\|^4}{4(d+2)^2}\,I.
\label{eq:Gt_squared}
\end{equation}

\emph{Leading-in-$d$ block.}
Conditional on $(\cH_{t-1},\theta_t)$, define $e_1 := \theta_t/\|\theta_t\|$
(any unit vector if $\theta_t=0$). This $e_1$ is now deterministic, so
the decomposition $u_t = z_t e_1 + u_{t,\perp}$ is well-defined with
$z_t := u_t^\top e_1\sim\cN(0,1)$ and $u_{t,\perp}$ independent of
$z_t$, distributed as $\cN(0,I-e_1e_1^\top)$ on the $(d{-}1)$-dim
subspace $e_1^\perp$. Crucially, $s_t = (z_t\|\theta_t\|+\varepsilon_t)^2-\hat\sigma^2$
depends on $u_t$ only through $z_t$, so $s_t$ is independent of
$u_{t,\perp}$ given $(\cH_{t-1},\theta_t)$.

Expand $u_tu_t^\top = z_t^2 e_1e_1^\top + z_t(e_1 u_{t,\perp}^\top + u_{t,\perp}e_1^\top) + u_{t,\perp}u_{t,\perp}^\top$
and $\|u_t\|^2 = z_t^2 + \|u_{t,\perp}\|^2$. The cross terms $z_t e_1 u_{t,\perp}^\top$
vanish under $\E[\cdot\mid\cH_{t-1},\theta_t]$ by symmetry of
$u_{t,\perp}$. The surviving contributions split into:
\begin{align*}
\E\bigl[s_t^2\|u_t\|^2\,z_t^2\,e_1e_1^\top \,\big|\,\cH_{t-1},\theta_t\bigr]
&= \Bigl(\E[s_t^2 z_t^4] + (d{-}1)\,\E[s_t^2 z_t^2]\Bigr)\,e_1e_1^\top, \\
\E\bigl[s_t^2\|u_t\|^2\,u_{t,\perp}u_{t,\perp}^\top \,\big|\,\cH_{t-1},\theta_t\bigr]
&= \E[s_t^2 z_t^2]\,(I-e_1e_1^\top) + \E[s_t^2]\,(d{+}1)(I-e_1e_1^\top),
\end{align*}
where conditional expectations are abbreviated, and we used:
\begin{itemize}[leftmargin=1.5em,nosep]
\item $\E[u_{t,\perp}u_{t,\perp}^\top]=I-e_1e_1^\top$ (covariance);
\item $\E[\|u_{t,\perp}\|^2]=d-1$ ($(d{-}1)$-dim chi-square);
\item $\E[\|u_{t,\perp}\|^2\,u_{t,\perp}u_{t,\perp}^\top]=(d+1)(I-e_1e_1^\top)$
  (Isserlis on the $(d{-}1)$-dim subspace, valid for $d\ge 2$);
\item independence of $u_{t,\perp}$ from $s_t$ given $(\cH_{t-1},\theta_t)$
  splits the joint expectation in the second display.
\end{itemize}

\emph{Bounding the conditional moments of $s_t$.} Write
$\xi_t := z_t\|\theta_t\|+\varepsilon_t$, so $s_t = \xi_t^2-\hat\sigma^2$ and
$s_t^2 = \xi_t^4 - 2\hat\sigma^2\xi_t^2 + \hat\sigma^4$. Given
$\theta_t$, $\xi_t$ has variance $\|\theta_t\|^2+\sigma_\varepsilon^2\le S_w^2+\sigma_\varepsilon^2$.
Using the assumed fourth-moment bound $\E[\varepsilon_t^4\mid\cH_{t-1}]\le 3\sigma_\varepsilon^4$
and Gaussian moments $\E[z_t^k]=k!!$ for $k\le 8$ (in particular
$\E[z_t^8]=105$ — needed because $s_t^2 z_t^4$ contains a $z_t^8$
contribution from $\xi_t^4 z_t^4 = (z_t\|\theta_t\|+\varepsilon_t)^4 z_t^4$),
straightforward expansion gives
\[
\E[s_t^2\mid\cH_{t-1},\theta_t]\le C_0\,D,\quad
\E[s_t^2 z_t^2\mid\cH_{t-1},\theta_t]\le C_0\,D,\quad
\E[s_t^2 z_t^4\mid\cH_{t-1},\theta_t]\le C_0\,D,
\]
with $D := S_w^4+\sigma_\varepsilon^4+\hat\sigma^4$ and a universal
constant $C_0$ (the dominant contribution is
$\E[\xi_t^4 z_t^4]\le 105\,S_w^4 + 90\,S_w^2\sigma_\varepsilon^2+9\sigma_\varepsilon^4
\le 320\,D$ via $\E[z_t^8]=105$ and AM-GM, so $C_0\le 320$ suffices).

\emph{Combining the blocks.} The parallel-block coefficient is
$A_1 := \E[s_t^2 z_t^4]+(d-1)\E[s_t^2 z_t^2]\le d\,C_0\,D$, and the
perpendicular-block coefficient is
$A_2 := \E[s_t^2 z_t^2]+(d+1)\E[s_t^2]\le (d+2)\,C_0\,D$. Hence
$\|\E[s_t^2\|u_t\|^2 u_tu_t^\top\mid\cH_{t-1},\theta_t]\|_\op
\le \max(A_1,A_2)\le (d+2)\,C_0\,D$. Multiplying by the
$G_t^2$-prefactor $d/(4(d+2))$ from \eqref{eq:Gt_squared} yields a
leading-block contribution
$\le (d/(4(d+2)))\cdot (d+2)\,C_0\,D = (C_0/4)\,d\,D$ exactly.

\emph{Trace correction (exact bound).} The second term of
\eqref{eq:Gt_squared} contributes
$\E[s_t^2\|u_t\|^4]/(4(d+2)^2)\cdot I$. Expanding
$\|u_t\|^4 = z_t^4 + 2z_t^2\|u_{t,\perp}\|^2 + \|u_{t,\perp}\|^4$ and
using $\E[\|u_{t,\perp}\|^4]=(d-1)(d+1)$,
\[
\E[s_t^2\|u_t\|^4\mid\cH_{t-1},\theta_t]
\;\le\;C_0\,D\,\bigl(1+2(d-1)+(d^2-1)\bigr)
\;=\;C_0\,D\,(d^2+2d-2).
\]
Divided by $4(d+2)^2$, the trace contribution to
$\|\E[G_t^2\mid\cH_{t-1},\theta_t]\|_\op$ is at most
$C_0\,D\,(d^2+2d-2)/(4(d+2)^2)$.

\emph{Conclusion (exact prefactor accounting).} Combining the
leading-block contribution $\le (C_0/4)\,d\,D$ with the trace
correction yields, after dividing by $d\,D$,
\[
\frac{\|\E[G_t^2\mid\cH_{t-1},\theta_t]\|_\op}{d\,D}
\;\le\;\frac{C_0}{4}
\;+\;\frac{C_0\,(d^2+2d-2)}{4d\,(d+2)^2}.
\]
With $C_0=320$, the second term is monotone decreasing in $d$ for
$d\ge 2$: at $d=2$, $C_0\cdot 6/(8\cdot 16)=15$; at $d=4$,
$C_0\cdot 22/(16\cdot 36)\approx 12.22$; as $d\to\infty$, the term
approaches $0$. Combined with the constant first term $C_0/4=80$,
the per-$dD$ bound is $\le 80+15=95$ at $d=2$, $\le 92.22$ at $d=4$,
and $\to 80$ as $d\to\infty$ — uniformly $\le 100$. Hence
$\sigma_*^2 = C\,D$ with $C = 100$ suffices for all $d\ge 2$.
Taking conditional expectation over $\theta_t$ preserves the bound
by linearity and operator-norm convexity, yielding
\eqref{eq:lifted_var}.
\end{proof}

\begin{theorem}[Variance-adaptive concentration; sharpening of Theorem~\ref{thm:matrix_bernstein_conf}]
\label{thm:matrix_bernstein_var}
Under the assumptions of Theorem~\ref{thm:matrix_bernstein_conf}
together with the additional hypotheses of Lemma~\ref{lem:lifted_variance}
($d\ge 2$ and $\E[\varepsilon_t^4\mid\cH_{t-1}]\le 3\sigma_\varepsilon^4$;
truncation by indicator $\mathbf 1\{\|u_t\|\le L\}$), with
probability at least $1-2\delta$ (jointly over the truncation event of
Lemma~\ref{lem:G_bound_conf} and the Bernstein event),
\begin{equation}
\bigl\|\widehat M_k-\bar M_k^{\mathrm{probe}}-\widetilde B\bigr\|_\op
\;\le\;\sigma_*\sqrt{\frac{2\,d\,\log(2d/\delta)}{m_k}}
\;+\;\frac{4 R_X\log(2d/\delta)}{3\,m_k}
\;+\;\|\Theta_k\|_\op,
\label{eq:bern_var}
\end{equation}
where $R_X = \Theta(L^2)=\Theta(d\log T)$ is the worst-case envelope of
Lemma~\ref{lem:G_bound_conf} and $\sigma_*^2$ is the variance constant
of Lemma~\ref{lem:lifted_variance}.
\end{theorem}
\begin{proof}
This is matrix Bernstein \citep[Theorem~7.7.1]{tropp2015introduction}
applied to the centered MDS $\{\tilde X_t\}_{t\in\cT_k}$, with:
\begin{itemize}[leftmargin=1.5em,nosep]
\item a.s.\ envelope $L_X := \|\tilde X_t\|_\op\le 2R_X$ (Lemma~\ref{lem:G_bound_conf}, plus centering);
\item predictable variance
$\bigl\|\sum_{t\in\cT_k}\E[\tilde X_t^2\mid\cH_{t-1}]\bigr\|_\op\le m_k\,\sigma_*^2\,d$
(Lemma~\ref{lem:lifted_variance}, summed).
\end{itemize}
Tropp's matrix Bernstein with these inputs gives, with probability
$\ge 1-\delta$,
\[
\Bigl\|\sum_{t\in\cT_k}\tilde X_t\Bigr\|_\op
\;\le\;\sqrt{2 m_k\sigma_*^2 d\log(2d/\delta)} + \frac{2L_X\log(2d/\delta)}{3}
\;\le\;\sqrt{2 m_k\sigma_*^2 d\log(2d/\delta)} + \frac{4 R_X\log(2d/\delta)}{3},
\]
using $L_X\le 2R_X$ in the linear-envelope term. Dividing by $m_k$
and combining with the truncation-bias offset $\|\Theta_k\|_\op$ via
the same transfer step as in Theorem~\ref{thm:matrix_bernstein_conf}
yields \eqref{eq:bern_var}. The two constituent events (truncation
and Bernstein) each have mass $\ge 1-\delta$; union bound gives
$1-2\delta$.
\end{proof}

\begin{corollary}[Sharpened projector concentration; sharpening of Corollary~\ref{cor:projector_conf}]
\label{cor:projector_var}
Under the assumptions of Corollary~\ref{cor:projector_conf} together
with the additional hypotheses of Lemma~\ref{lem:lifted_variance}
($d\ge 2$, $\E[\varepsilon_t^4\mid\cH_{t-1}]\le 3\sigma_\varepsilon^4$,
indicator-truncated probes), with probability at least $1-2\delta$,
\begin{equation}
\varepsilon_k = \|\widehat P_k - P_k^\star\|_\op
\;\le\; C_{\mathrm{sub}}^\star\,\sqrt{\frac{d\,\log(2d/\delta)}{m_k}}
\;+\;\frac{16\,R_X\,\log(2d/\delta)}{3\,\lambda_{\min}\,m_k}
\;+\;\Delta_\sigma\;+\;\frac{4\|\Theta_k\|_\op}{\lambda_{\min}},
\quad
C_{\mathrm{sub}}^\star \;:=\;\frac{4\sqrt{2}\,\sigma_*}{\lambda_{\min}}.
\label{eq:proj_bound_var}
\end{equation}
The linear-in-$1/m_k$ term is dominated by the leading
$\sqrt{d/m_k}$ term whenever $m_k\gtrsim R_X^2\log(2d/\delta)/(\sigma_*^2 d)$,
which evaluates to $m_k\gtrsim d\log^2 T\cdot\log(2d/\delta) = \tilde\Theta(d)$
(with $R_X=\Theta(d\log T)$ from Lemma~\ref{lem:G_bound_conf}). The
$\tilde\Theta$ hides polylog factors of $T$ and $d/\delta$.
The sharpened constant $C_{\mathrm{sub}}^\star = \Theta(1/\lambda_{\min})$
(in the variance scale) gives leading-term scaling
$O(\sqrt{d/m_k})$, replacing the original
$C_{\mathrm{sub}}=\Theta(d)/\lambda_{\min}$ leading scaling
$O(d/\sqrt{m_k})$ — a factor-of-$\sqrt d$ improvement in the
$m_k$-dependence at any fixed $d$.
\end{corollary}
\begin{proof}
Davis--Kahan applied to \eqref{eq:bern_var}, with the leading
$\sigma_*\sqrt{2d\log(2d/\delta)/m_k}$ replacing the original
$2R_X\sqrt{\log(2d/\delta)/m_k}=\Theta(d\log T\sqrt{\log(2d/\delta)/m_k})$.
The other terms (linear $1/m_k$ Bernstein term, truncation bias,
variance-misspecification floor $\Delta_\sigma$) are unchanged.
\end{proof}

\begin{theorem}[Sharpened probe term; sharpening of Theorem~\ref{thm:spsc_regret}]
\label{thm:spsc_regret_sharp}
In addition to the assumptions of Theorem~\ref{thm:spsc_regret},
assume the conditions of Lemma~\ref{lem:lifted_variance}: $d\ge 2$
and $\E[\varepsilon_t^4\mid\cH_{t-1}]\le 3\sigma_\varepsilon^4$
(satisfied by Gaussian noise; truncation is by indicator
$\mathbf 1\{\|u_t\|\le L\}$, so $\mathcal K^{-1}$ is the full-Gaussian
inverse used throughout). Run Alg.~\ref{alg:spsc} with the sharpened
probe allocation $m_k = \lceil(B^\star\ell_k/(2A))^{2/3}\rceil$, where
\[
A := 2R_\cA S_w + c, \qquad
B^\star := 4 C_{\mathrm{sub}}^\star\,S_w\,R_\cA\,\sqrt{d\,\log(2d/\delta)}
\;=\;\Theta(\sqrt d).
\]
Assume the per-segment regime
$\ell_k\ge \ell^\star = \tilde\Theta(d)$ for all $k\in[K]$
(explicitly, $\ell^\star = C_\ell\,d\,\mathrm{polylog}(T,d/\delta)$
for a universal constant $C_\ell$, where the polylog absorbs the
$\log(2d/\delta)$ and $\log T$ factors of Corollary~\ref{cor:projector_var};
when segments are balanced, this is equivalent to
$T\ge K\ell^\star = \tilde\Theta(Kd)$).
Then, with probability $\ge 1-\delta$,
\begin{equation}
\DynReg_T^{(c)}
\;=\;\tilO\!\bigl(r\sqrt{KT}\bigr)
\;+\;\tilO\!\bigl(d^{1/3}\,K^{1/3}\,T^{2/3}\bigr)
\;+\;O(WV)\;+\;O(T\Delta_\sigma).
\label{eq:main_bound_sharp}
\end{equation}
The probe-term prefactor improves from $d^{2/3}$ to $d^{1/3}$.
\end{theorem}
\begin{proof}
The proof of Theorem~\ref{thm:spsc_regret} (\S\ref{app:proof_main})
uses $\varepsilon_k\le C_{\mathrm{sub}}\sqrt{\log(2d/\delta)/m_k}+\Delta_\sigma$
inside the step (i.c) tradeoff
\[
2R_\cA S_w\,\varepsilon_k\,n_k \;\le\; B\,n_k/\sqrt{m_k} + 2R_\cA S_w\Delta_\sigma\,n_k,
\quad B = 4 C_{\mathrm{sub}}\,S_w\,R_\cA\sqrt{\log(2d/\delta)}.
\]
Corollary~\ref{cor:projector_var} replaces this with
\[
\varepsilon_k \le C_{\mathrm{sub}}^\star\sqrt{\frac{d\log(2d/\delta)}{m_k}}
+\frac{16 R_X\log(2d/\delta)}{3\lambda_{\min}m_k}+\Delta_\sigma.
\]
Multiplying by $2R_\cA S_w n_k$, the leading variance term becomes
$B^\star n_k/\sqrt{m_k}$ with $B^\star = \Theta(\sqrt d)$ as displayed,
and the linear-Bernstein term contributes
$32 R_\cA S_w R_X \log(2d/\delta)\,n_k/(3\lambda_{\min}m_k) = \tilO(d\,n_k/m_k)$.

\emph{Threshold derivation.} From Corollary~\ref{cor:projector_var},
the linear-Bernstein contribution to the regret tradeoff is
$\tilO(d\,n_k/m_k)$ and the leading variance contribution is
$\tilO(\sqrt d\,n_k/\sqrt{m_k})$ (both up to polylog factors of
$T$ and $d/\delta$, which differ between the two terms but get
absorbed into $\tilde\Theta$ uniformly below). Their ratio at the sharpened optimum
$m_k^\star\asymp d^{1/3}\ell_k^{2/3}$ is
\[
\frac{\text{leading}}{\text{linear}}
\;\asymp\;
\frac{\sqrt{d/m_k^\star}}{d/m_k^\star}
\;=\;\sqrt{m_k^\star/d}
\;\asymp\;(\ell_k/d)^{1/3}.
\]
For domination $\ge 1$ (modulo polylog), $\ell_k\gtrsim d$ up to
polylog factors. We absorb these into $\ell^\star = \tilde\Theta(d)$;
the explicit form $C_\ell\,d\,\mathrm{polylog}(T,d/\delta)$ handles
both the $R_X = O(d\log T)$ envelope and the $\log(2d/\delta)$
Bernstein-tail factor.

\emph{Per-segment optimization.} In the regime $\ell_k\ge\ell^\star$,
the linear term is dominated, and
$\min_{m_k}\bigl(A m_k+B^\star\ell_k/\sqrt{m_k}\bigr)$
is achieved at $m_k^\star = (B^\star\ell_k/(2A))^{2/3}$ with value
\[
\Theta\bigl(A^{1/3}(B^\star)^{2/3}\bigr)\,\ell_k^{2/3} = \Theta(d^{1/3})\,\ell_k^{2/3}
\]
per segment. The leading-to-linear ratio at $m_k^\star$ is
$(B^\star\ell_k/\sqrt{m_k^\star})/(d\,n_k/m_k^\star)
\asymp (\ell_k/d)^{1/3}\to\infty$ under $\ell_k\ge\ell^\star$, confirming
domination.

\emph{Summing over segments.} Under the theorem's hypothesis
$\ell_k\ge\ell^\star$ for all $k\in[K]$, the per-segment
contributions sum via Jensen
$\sum_k\ell_k^{2/3}\le K^{1/3}T^{2/3}$ to
$\tilO(d^{1/3}K^{1/3}T^{2/3})$. (For balanced segments
$\ell_k=T/K$, the hypothesis reduces to $T\ge K\ell^\star=\tilde\Theta(Kd)$,
satisfied for $T$ large at any fixed $K,d$.)

All other steps of the proof of Theorem~\ref{thm:spsc_regret} (the
windowed-ridge in-subspace argument, the drift summation, the
variance-misspecification floor) carry over verbatim because they do
not involve $C_{\mathrm{sub}}$.\qed
\end{proof}

\begin{remark}[Algorithm structure unchanged; recommended probe rate updated]
\label{rem:no_algo_change}
Theorem~\ref{thm:spsc_regret_sharp} preserves the algorithmic
structure of Alg.~\ref{alg:spsc} (probe-then-exploit interleaving,
projector recovery via $\widehat M_k$, windowed projected ridge-UCB);
only the recommended probe rate is updated:
$m_k^\star\asymp d^{1/3}\ell_k^{2/3}$ replaces the original
$m_k^\star\asymp d^{2/3}\ell_k^{2/3}$, a $d^{1/3}$-factor reduction.
The empirical results of \S\ref{sec:experiments} were obtained with
the original schedule, so they correspond to a (theoretically)
suboptimal — but consistent — probe rate; reported regret is therefore
an upper bound on what the sharpened schedule would achieve.
The empirical $d$-scaling slope $\approx-0.16$ in Figure~\ref{fig:d-scaling}
is closer to the sharpened $d^{1/3}=+0.33$ prediction than to the
original $d^{2/3}=+0.67$, providing post-hoc empirical support for
the variance-adaptive analysis.
\end{remark}

\begin{remark}[Combined LB/UB picture, after both extensions]
\label{rem:combined_after_ext}
Theorem~\ref{thm:spsc_lb_piecewise} and
Theorem~\ref{thm:spsc_regret_sharp} together yield
\[
\Omega\!\bigl(c^{1/3}\,K^{1/3}T^{2/3}\bigr)
\;\le\;\text{(probe-acquisition regret of SPSC)}\;\le\;
\tilO\!\bigl(d^{1/3}\,K^{1/3}T^{2/3}\bigr),
\]
tight in $K$ and $T$, with a residual $d^{1/3}$ gap that the
empirical $d$-scaling fit (Figure~\ref{fig:d-scaling}) suggests is
also loose; closing it likely requires the subspace-restricted
concentration sketched in Remark~\ref{rem:dim_dep}, which is beyond
the scope of this purely-additive sharpening.
\end{remark}

% =====================================================================
% End extensions.tex
% =====================================================================
` to conf.tex just before `\end{document}`
%     (i.e., after Section J / Table tab:published_comparison).
%   * Optionally add an entry to the appendix TOC (line ~847 of conf.tex):
%       \tocentry{K}{app:extensions}{Extensions: piecewise LB and sharpened d-dependence}
% =====================================================================

\section{Extensions: piecewise lower bound and sharpened \texorpdfstring{$d$}{d}-dependence}
\label{app:extensions}

This section supplements (without modifying) Theorems~\ref{thm:spsc_regret}--\ref{thm:spsc_lb}
with two additive results: a piecewise-stationary lower bound that
matches the upper-bound rate in $K$ (\S\ref{app:ext_piecewise_lb}),
and a sharpened $d$-dependence in the probe term via variance-adaptive
matrix Bernstein (\S\ref{app:ext_sharpened_d}). The original theorems
remain in force; the results below are stated as \emph{independent}
sharpenings.

\subsection{Piecewise-stationary lower bound}
\label{app:ext_piecewise_lb}

Theorem~\ref{thm:spsc_lb} establishes the LB only for $K=1$ (single
segment of length $T$). The $K$-segment extension follows by applying
the single-segment bound conditionally inside each segment, exploiting
the fact that the per-segment hidden direction $v_k$ in the
construction is independent of all previous segments. No new technical
tools beyond Theorem~\ref{thm:spsc_lb} itself are required.

\begin{theorem}[Piecewise lower bound; extension of Theorem~\ref{thm:spsc_lb} to $K\ge 1$]
\label{thm:spsc_lb_piecewise}
Let $\Pi$ be the policy class of Theorem~\ref{thm:spsc_lb} (bounded
exploit actions $\|x_t\|\le R_\cA$, isotropic Gaussian probes
$u_t\sim\cN(0,I_d)$). Fix $K\ge 1$ and segment lengths
$\{\ell_k\}_{k=1}^K$ with $\sum_k\ell_k=T$. There exist constants
$c_0,d_0,\ell_0>0$ depending only on $R_\cA,\sigma_\varepsilon,c$
(in fact, the same constants as Theorem~\ref{thm:spsc_lb}) such that
whenever $\min_k\ell_k\ge\ell_0$ and $d\ge d_0\,(\min_k\ell_k)^{1/3}$,
\begin{equation}
\inf_{\pi\in\Pi}\sup_\nu\ \E^\pi_\nu[\DynReg_T^{(c)}]
\;\ge\; c_0\,c^{1/3}(R_\cA\sigma_\varepsilon)^{2/3}\sum_{k=1}^K\ell_k^{2/3}.
\label{eq:lb_piecewise}
\end{equation}
For uniform segments $\ell_k=T/K$ this yields
\begin{equation}
\inf_{\pi\in\Pi}\sup_\nu\ \E^\pi_\nu[\DynReg_T^{(c)}]
\;\ge\; c_0\,c^{1/3}(R_\cA\sigma_\varepsilon)^{2/3}\,K^{1/3}T^{2/3},
\label{eq:lb_piecewise_uniform}
\end{equation}
matching the probe-acquisition rate of Theorem~\ref{thm:spsc_regret}
in both $K$ and $T$ (with the same residual $d$-prefactor gap discussed
in Remark~\ref{rem:dim_dep}, sharpened in
\S\ref{app:ext_sharpened_d} below).
\end{theorem}

\begin{proof}
\emph{Construction.} Let $v_1,\dots,v_K$ be i.i.d.\ uniform on
$\mathbb S^{d-1}$, drawn independently of the policy. For
per-segment scales $\alpha_k>0$ to be specified, define the instance
$\nu(\vec v):\theta_t = \alpha_k v_k$ for $t\in\cI_k$. Let
$\nu_0:\theta_t\equiv 0$ denote the global null. The per-segment
``zero-on-segment-$k$'' alternative is
$\nu^{(-k)}(\vec v):\theta_t = \alpha_{k'}v_{k'}\mathbf 1\{k'\ne k\}$
for $t\in\cI_{k'}$ — i.e., identical to $\nu(\vec v)$ except on
segment $k$ where $\theta_t=0$.

\emph{Per-segment KL.} The trajectory laws under $\nu(\vec v)$ and
$\nu^{(-k)}(\vec v)$ agree on segments $k'\ne k$ and differ only on
segment $k$. The sequential KL chain rule therefore gives
\[
\mathrm{KL}\bigl(\mathbb P_{\nu^{(-k)}(\vec v)}\,\bigl\|\,\mathbb P_{\nu(\vec v)}\bigr)
\;=\; \frac{\alpha_k^2}{2\sigma_\varepsilon^2}\,
\E^\pi_{\nu^{(-k)}(\vec v)}\!\Bigl[\sum_{t\in\cI_k}(x_t^\top v_k)^2\Bigr].
\]
Under $\nu^{(-k)}(\vec v)$, on segment $k$ the parameter is $0$, so
the segment-$k$ marginal trajectory is independent of $v_k$.
Consequently $v_k$ is independent of the actions $\{x_t\}_{t\in\cI_k}$
under $\nu^{(-k)}(\vec v)$, and the trace-identity argument of
Lemma~\ref{lem:trace_kl_conf} applies verbatim with horizon $\ell_k$,
probe count $m_k$, and exploit count $n_k=\ell_k-m_k$:
\[
\E_{v_k}\E^\pi_{\nu^{(-k)}(\vec v)}\!\Bigl[\sum_{t\in\cI_k}(x_t^\top v_k)^2\Bigr]
\;\le\; m_k + R_\cA^2 n_k/d \;=:\;Q_k,
\]
where the average is over $v_k\sim\mathrm{Unif}(\mathbb S^{d-1})$. As in
Lemma~\ref{lem:trace_kl_conf}, the existence of an extremal direction
gives a deterministic $v_k^\star\in\mathbb S^{d-1}$ for which the
inequality holds with $v_k^\star$ replacing the average. Setting
$\alpha_k^2 := \sigma_\varepsilon^2/(4 Q_k)$ yields
\[
\mathrm{KL}\bigl(\mathbb P_{\nu^{(-k)}(\vec v^\star)}\,\bigl\|\,\mathbb P_{\nu(\vec v^\star)}\bigr)
\;\le\;\frac{\alpha_k^2 Q_k}{2\sigma_\varepsilon^2}\;=\;\frac{1}{8},
\]
hence Pinsker gives per-segment TV $\le 1/4$.

\emph{Per-segment regret transfer.} The optimal action under
$\nu(\vec v^\star)$ on segment $k$ is $x_t^\star=R_\cA v_k^\star$ with
reward $\alpha_k R_\cA$; under $\nu^{(-k)}(\vec v^\star)$ the
parameter on segment $k$ is $0$ so any action is optimal, and the
segment-$k$ regret of the policy is identically $0$. The
arm-regret-vs.-alignment manipulation of \S\ref{app:proof_lb}
(equation~\eqref{eq:reg_to_alignment_conf}), applied to the segment-$k$
exploit rounds, gives
\[
\E^\pi_{\nu(\vec v^\star)}\!\Bigl[\sum_{t\in\cI_k\cap\cT_{\mathrm{ex}}}\Delta_t\Bigr]
\;\ge\;\alpha_k n_k R_\cA\;-\;\alpha_k\,\E^\pi_{\nu(\vec v^\star)}[f_k],\quad
f_k:=\sum_{t\in\cI_k\cap\cT_{\mathrm{ex}}}|v_k^{\star\top}x_t|.
\]
TV transfer from $\nu(\vec v^\star)$ to $\nu^{(-k)}(\vec v^\star)$
with $\|f_k\|_\infty\le n_k R_\cA$ and segment-$k$
$d_{\mathrm{TV}}\le 1/4$:
\[
\E^\pi_{\nu(\vec v^\star)}[f_k]\;\le\;\E^\pi_{\nu^{(-k)}(\vec v^\star)}[f_k]+\tfrac12\,n_k R_\cA.
\]
Under $\nu^{(-k)}(\vec v^\star)$, by Cauchy--Schwarz and the
trace-identity bound applied on segment $k$,
$\E^\pi_{\nu^{(-k)}(\vec v^\star)}[f_k]\le R_\cA n_k/\sqrt d$.
Combining (with $d\ge 16$ so $1/\sqrt d\le 1/4$):
\begin{equation}
\E^\pi_{\nu(\vec v^\star)}\!\Bigl[\sum_{t\in\cI_k\cap\cT_{\mathrm{ex}}}\Delta_t\Bigr]
\;\ge\;\tfrac14\,\alpha_k n_k R_\cA
\;=\;\frac{R_\cA\sigma_\varepsilon\,n_k}{8\sqrt{Q_k}}.
\label{eq:per_seg_regret_lb}
\end{equation}

\emph{Per-segment costed regret.} Adding the segment-$k$ probe cost
$cm_k$,
\[
\E^\pi_{\nu(\vec v^\star)}[\DynReg^{(c)}_k]
\;\ge\; cm_k + \frac{R_\cA\sigma_\varepsilon\,n_k}{8\sqrt{Q_k}},
\]
with the same functional form as the single-segment bound
\eqref{eq:dynreg_lower_conf}. The Regime-I/Regime-II optimization of
\S\ref{app:proof_lb} applies verbatim with $T\to\ell_k$, yielding
\[
\E^\pi_{\nu(\vec v^\star)}[\DynReg^{(c)}_k]
\;\ge\;c_0\,c^{1/3}(R_\cA\sigma_\varepsilon)^{2/3}\ell_k^{2/3}
\]
in the regime $\ell_k\ge\ell_0$ and $d\ge d_0\,\ell_k^{1/3}$, with
$c_0,d_0,\ell_0$ as in Theorem~\ref{thm:spsc_lb}.

\emph{Summing across segments.} Total dynamic regret is the sum of
per-segment costed regrets (probes and exploit rounds partition by
segment):
\[
\E^\pi_{\nu(\vec v^\star)}[\DynReg^{(c)}_T]
\;=\;\sum_{k=1}^K\E^\pi_{\nu(\vec v^\star)}[\DynReg^{(c)}_k]
\;\ge\;c_0\,c^{1/3}(R_\cA\sigma_\varepsilon)^{2/3}\sum_{k=1}^K\ell_k^{2/3}.
\]
Taking $\sup_\nu\ge$ this specific instance yields
\eqref{eq:lb_piecewise}; specialising to $\ell_k=T/K$ gives
$\sum_k\ell_k^{2/3}=K\,(T/K)^{2/3}=K^{1/3}T^{2/3}$, hence
\eqref{eq:lb_piecewise_uniform}.\qed
\end{proof}

\begin{remark}[The argument uses only Theorem~\ref{thm:spsc_lb}'s tools]
\label{rem:lb_piecewise_tools}
The proof reuses the trace-identity Lemma~\ref{lem:trace_kl_conf} and
the Le~Cam transfer of \S\ref{app:proof_lb} per segment. The only new
ingredient is the observation that the segment-$k$ KL between
$\nu(\vec v^\star)$ and $\nu^{(-k)}(\vec v^\star)$ is exactly the
single-segment KL with parameters $(\ell_k, m_k, n_k)$, because the
two laws agree outside segment $k$. No multi-hypothesis Fano, no
sliding-window reduction, and no new concentration is invoked.
\end{remark}

\begin{remark}[Adaptive probe schedules]
\label{rem:lb_piecewise_adaptive}
The bound \eqref{eq:per_seg_regret_lb} is in terms of the policy's
\emph{realised} segment-$k$ probe count $m_k$ and exploit count $n_k$,
which may be adaptive (chosen based on history through $\tau_{k-1}$).
The per-segment Regime-I/Regime-II optimization minimizes
$cm_k+R_\cA\sigma_\varepsilon n_k/(8\sqrt{Q_k})$ over $m_k\in[0,\ell_k]$;
the policy's adaptive choice of $m_k$ can only weakly increase this
quantity. Hence the LB applies under any (possibly adaptive)
per-segment probe allocation.
\end{remark}

\begin{remark}[Combined LB/UB picture]
\label{rem:combined_lb_ub}
Theorem~\ref{thm:spsc_lb_piecewise} combined with
Theorem~\ref{thm:spsc_regret} yields the matching pair
\[
\Omega(c^{1/3}\,K^{1/3}T^{2/3})
\;\le\;\text{(probe-acquisition regret)}\;\le\;
\tilO(d^{2/3}\,K^{1/3}T^{2/3}),
\]
tight in $K$ and $T$. The remaining gap is the $d$-prefactor on the
upper-bound side; \S\ref{app:ext_sharpened_d} narrows that gap from
$d^{2/3}$ to $d^{1/3}$ via a sharper concentration argument, leaving
a residual $d^{1/3}$ open question.
\end{remark}

\subsection{Sharpened \texorpdfstring{$d$}{d}-dependence via variance-adaptive Bernstein}
\label{app:ext_sharpened_d}

Theorem~\ref{thm:matrix_bernstein_conf} (matrix Freedman) bounds
$\|\widehat M_k-\bar M_k^{\mathrm{probe}}-\widetilde B\|_\op$ using the
worst-case envelope $\|\tilde X_t\|_\op\le 2R_X$ with
$R_X=\Theta(L^2)=\Theta(d\log T)$. The leading term
$2R_X\sqrt{\log(2d/\delta)/m_k}$ then yields
$C_{\mathrm{sub}}=8R_X/\lambda_{\min}=\Theta(d\log T/\lambda_{\min})$
in Corollary~\ref{cor:projector_conf}, and the constant
$B = 4 C_{\mathrm{sub}} S_w R_\cA\sqrt{\log(2d/\delta)} = \Theta(d)$
propagates through the proof of Theorem~\ref{thm:spsc_regret} to give
the $\tilO(d^{2/3}K^{1/3}T^{2/3})$ probe-acquisition rate.

The envelope is loose. For isotropic Gaussian $u_t$,
$\|u_t\|^2$ concentrates near $d$ with $\sqrt d$ fluctuations, so the
\emph{predictable second moment} of $\tilde G_t$ in operator norm is
$\Theta(d)$ per round, not $\Theta(d^2)$ as the envelope suggests.
The variance-adaptive matrix Bernstein
inequality~\citep[Theorem~7.7.1]{tropp2015introduction} exploits this
gap and yields a sharper concentration. The algorithmic structure of
SPSC (probe-then-exploit interleaving, projector recovery, projected
ridge-UCB) is unchanged; only the recommended probe rate $m_k$ is
updated by a $d^{1/3}$ factor
(from $m_k^\star\asymp d^{2/3}\ell_k^{2/3}$ to
$m_k^\star\asymp d^{1/3}\ell_k^{2/3}$; see
Theorem~\ref{thm:spsc_regret_sharp} and Remark~\ref{rem:no_algo_change}).

\begin{lemma}[Predictable second moment of the lifted probe sample]
\label{lem:lifted_variance}
Assume $d\ge 2$, $\|\theta_t\|\le S_w$ a.s., and that the noise has a
bounded fourth moment $\E[\varepsilon_t^4\mid\cH_{t-1}]\le 3\sigma_\varepsilon^4$
(satisfied by Gaussian $\varepsilon_t\sim\cN(0,\sigma_\varepsilon^2)$).
Let $u_t\sim\cN(0,I_d)$ on probe rounds (truncation by an indicator
$\mathbf 1\{\|u_t\|\le L\}$ only decreases the second moment in the
PSD order, so the bound below applies; if probes are instead
distributed $u_t\mid\{\|u_t\|\le L\}$, the same bound holds up to a
$1+o(1)$ factor as $L\to\infty$). Define $G_t := \mathcal K^{-1}(s_t u_t u_t^\top)$
and $\tilde X_t := G_t\mathbf 1\{\cA_t\}-\E[G_t\mathbf 1\{\cA_t\}\mid\cH_{t-1}]$.
Then there exists $\sigma_*^2$ depending only on $S_w,\sigma_\varepsilon,\hat\sigma$
(and not on $d$) such that
\begin{equation}
\bigl\|\E[\tilde X_t^2\mid\cH_{t-1}]\bigr\|_\op
\;\le\;\sigma_*^2\,d.
\label{eq:lifted_var}
\end{equation}
Explicitly, $\sigma_*^2 \le C\,(S_w^4 + \sigma_\varepsilon^4 + \hat\sigma^4)$
with $C=100$ (using the exact trace-expansion bound established in the
proof; the value of $C$ does not affect the leading-order rate).
\end{lemma}
\begin{proof}
Centering only reduces the second moment in the PSD order:
$\E[\tilde X_t^2\mid\cH_{t-1}]\preceq\E[G_t^2\mathbf 1\{\cA_t\}\mid\cH_{t-1}]\preceq\E[G_t^2\mid\cH_{t-1}]$.
We bound $\|\E[G_t^2\mid\cH_{t-1}]\|_\op$ by conditioning further on
$\theta_t$ and proving a bound uniform in $\theta_t$ with
$\|\theta_t\|\le S_w$. By tower-property the same bound then applies
to $\E[\,\cdot\mid\cH_{t-1}]$.

By Lemma~\ref{lem:K_inverse},
\(
G_t = \tfrac{s_t}{2}\,u_t u_t^\top - \tfrac{s_t\|u_t\|^2}{2(d{+}2)}\,I.
\)
Squaring and using $(uu^\top)^2 = \|u\|^2 uu^\top$:
\begin{equation}
G_t^2 = \frac{d}{4(d+2)}\,s_t^2\|u_t\|^2\, u_t u_t^\top
       \;+\;\frac{s_t^2\,\|u_t\|^4}{4(d+2)^2}\,I.
\label{eq:Gt_squared}
\end{equation}

\emph{Leading-in-$d$ block.}
Conditional on $(\cH_{t-1},\theta_t)$, define $e_1 := \theta_t/\|\theta_t\|$
(any unit vector if $\theta_t=0$). This $e_1$ is now deterministic, so
the decomposition $u_t = z_t e_1 + u_{t,\perp}$ is well-defined with
$z_t := u_t^\top e_1\sim\cN(0,1)$ and $u_{t,\perp}$ independent of
$z_t$, distributed as $\cN(0,I-e_1e_1^\top)$ on the $(d{-}1)$-dim
subspace $e_1^\perp$. Crucially, $s_t = (z_t\|\theta_t\|+\varepsilon_t)^2-\hat\sigma^2$
depends on $u_t$ only through $z_t$, so $s_t$ is independent of
$u_{t,\perp}$ given $(\cH_{t-1},\theta_t)$.

Expand $u_tu_t^\top = z_t^2 e_1e_1^\top + z_t(e_1 u_{t,\perp}^\top + u_{t,\perp}e_1^\top) + u_{t,\perp}u_{t,\perp}^\top$
and $\|u_t\|^2 = z_t^2 + \|u_{t,\perp}\|^2$. The cross terms $z_t e_1 u_{t,\perp}^\top$
vanish under $\E[\cdot\mid\cH_{t-1},\theta_t]$ by symmetry of
$u_{t,\perp}$. The surviving contributions split into:
\begin{align*}
\E\bigl[s_t^2\|u_t\|^2\,z_t^2\,e_1e_1^\top \,\big|\,\cH_{t-1},\theta_t\bigr]
&= \Bigl(\E[s_t^2 z_t^4] + (d{-}1)\,\E[s_t^2 z_t^2]\Bigr)\,e_1e_1^\top, \\
\E\bigl[s_t^2\|u_t\|^2\,u_{t,\perp}u_{t,\perp}^\top \,\big|\,\cH_{t-1},\theta_t\bigr]
&= \E[s_t^2 z_t^2]\,(I-e_1e_1^\top) + \E[s_t^2]\,(d{+}1)(I-e_1e_1^\top),
\end{align*}
where conditional expectations are abbreviated, and we used:
\begin{itemize}[leftmargin=1.5em,nosep]
\item $\E[u_{t,\perp}u_{t,\perp}^\top]=I-e_1e_1^\top$ (covariance);
\item $\E[\|u_{t,\perp}\|^2]=d-1$ ($(d{-}1)$-dim chi-square);
\item $\E[\|u_{t,\perp}\|^2\,u_{t,\perp}u_{t,\perp}^\top]=(d+1)(I-e_1e_1^\top)$
  (Isserlis on the $(d{-}1)$-dim subspace, valid for $d\ge 2$);
\item independence of $u_{t,\perp}$ from $s_t$ given $(\cH_{t-1},\theta_t)$
  splits the joint expectation in the second display.
\end{itemize}

\emph{Bounding the conditional moments of $s_t$.} Write
$\xi_t := z_t\|\theta_t\|+\varepsilon_t$, so $s_t = \xi_t^2-\hat\sigma^2$ and
$s_t^2 = \xi_t^4 - 2\hat\sigma^2\xi_t^2 + \hat\sigma^4$. Given
$\theta_t$, $\xi_t$ has variance $\|\theta_t\|^2+\sigma_\varepsilon^2\le S_w^2+\sigma_\varepsilon^2$.
Using the assumed fourth-moment bound $\E[\varepsilon_t^4\mid\cH_{t-1}]\le 3\sigma_\varepsilon^4$
and Gaussian moments $\E[z_t^k]=k!!$ for $k\le 8$ (in particular
$\E[z_t^8]=105$ — needed because $s_t^2 z_t^4$ contains a $z_t^8$
contribution from $\xi_t^4 z_t^4 = (z_t\|\theta_t\|+\varepsilon_t)^4 z_t^4$),
straightforward expansion gives
\[
\E[s_t^2\mid\cH_{t-1},\theta_t]\le C_0\,D,\quad
\E[s_t^2 z_t^2\mid\cH_{t-1},\theta_t]\le C_0\,D,\quad
\E[s_t^2 z_t^4\mid\cH_{t-1},\theta_t]\le C_0\,D,
\]
with $D := S_w^4+\sigma_\varepsilon^4+\hat\sigma^4$ and a universal
constant $C_0$ (the dominant contribution is
$\E[\xi_t^4 z_t^4]\le 105\,S_w^4 + 90\,S_w^2\sigma_\varepsilon^2+9\sigma_\varepsilon^4
\le 320\,D$ via $\E[z_t^8]=105$ and AM-GM, so $C_0\le 320$ suffices).

\emph{Combining the blocks.} The parallel-block coefficient is
$A_1 := \E[s_t^2 z_t^4]+(d-1)\E[s_t^2 z_t^2]\le d\,C_0\,D$, and the
perpendicular-block coefficient is
$A_2 := \E[s_t^2 z_t^2]+(d+1)\E[s_t^2]\le (d+2)\,C_0\,D$. Hence
$\|\E[s_t^2\|u_t\|^2 u_tu_t^\top\mid\cH_{t-1},\theta_t]\|_\op
\le \max(A_1,A_2)\le (d+2)\,C_0\,D$. Multiplying by the
$G_t^2$-prefactor $d/(4(d+2))$ from \eqref{eq:Gt_squared} yields a
leading-block contribution
$\le (d/(4(d+2)))\cdot (d+2)\,C_0\,D = (C_0/4)\,d\,D$ exactly.

\emph{Trace correction (exact bound).} The second term of
\eqref{eq:Gt_squared} contributes
$\E[s_t^2\|u_t\|^4]/(4(d+2)^2)\cdot I$. Expanding
$\|u_t\|^4 = z_t^4 + 2z_t^2\|u_{t,\perp}\|^2 + \|u_{t,\perp}\|^4$ and
using $\E[\|u_{t,\perp}\|^4]=(d-1)(d+1)$,
\[
\E[s_t^2\|u_t\|^4\mid\cH_{t-1},\theta_t]
\;\le\;C_0\,D\,\bigl(1+2(d-1)+(d^2-1)\bigr)
\;=\;C_0\,D\,(d^2+2d-2).
\]
Divided by $4(d+2)^2$, the trace contribution to
$\|\E[G_t^2\mid\cH_{t-1},\theta_t]\|_\op$ is at most
$C_0\,D\,(d^2+2d-2)/(4(d+2)^2)$.

\emph{Conclusion (exact prefactor accounting).} Combining the
leading-block contribution $\le (C_0/4)\,d\,D$ with the trace
correction yields, after dividing by $d\,D$,
\[
\frac{\|\E[G_t^2\mid\cH_{t-1},\theta_t]\|_\op}{d\,D}
\;\le\;\frac{C_0}{4}
\;+\;\frac{C_0\,(d^2+2d-2)}{4d\,(d+2)^2}.
\]
With $C_0=320$, the second term is monotone decreasing in $d$ for
$d\ge 2$: at $d=2$, $C_0\cdot 6/(8\cdot 16)=15$; at $d=4$,
$C_0\cdot 22/(16\cdot 36)\approx 12.22$; as $d\to\infty$, the term
approaches $0$. Combined with the constant first term $C_0/4=80$,
the per-$dD$ bound is $\le 80+15=95$ at $d=2$, $\le 92.22$ at $d=4$,
and $\to 80$ as $d\to\infty$ — uniformly $\le 100$. Hence
$\sigma_*^2 = C\,D$ with $C = 100$ suffices for all $d\ge 2$.
Taking conditional expectation over $\theta_t$ preserves the bound
by linearity and operator-norm convexity, yielding
\eqref{eq:lifted_var}.
\end{proof}

\begin{theorem}[Variance-adaptive concentration; sharpening of Theorem~\ref{thm:matrix_bernstein_conf}]
\label{thm:matrix_bernstein_var}
Under the assumptions of Theorem~\ref{thm:matrix_bernstein_conf}
together with the additional hypotheses of Lemma~\ref{lem:lifted_variance}
($d\ge 2$ and $\E[\varepsilon_t^4\mid\cH_{t-1}]\le 3\sigma_\varepsilon^4$;
truncation by indicator $\mathbf 1\{\|u_t\|\le L\}$), with
probability at least $1-2\delta$ (jointly over the truncation event of
Lemma~\ref{lem:G_bound_conf} and the Bernstein event),
\begin{equation}
\bigl\|\widehat M_k-\bar M_k^{\mathrm{probe}}-\widetilde B\bigr\|_\op
\;\le\;\sigma_*\sqrt{\frac{2\,d\,\log(2d/\delta)}{m_k}}
\;+\;\frac{4 R_X\log(2d/\delta)}{3\,m_k}
\;+\;\|\Theta_k\|_\op,
\label{eq:bern_var}
\end{equation}
where $R_X = \Theta(L^2)=\Theta(d\log T)$ is the worst-case envelope of
Lemma~\ref{lem:G_bound_conf} and $\sigma_*^2$ is the variance constant
of Lemma~\ref{lem:lifted_variance}.
\end{theorem}
\begin{proof}
This is matrix Bernstein \citep[Theorem~7.7.1]{tropp2015introduction}
applied to the centered MDS $\{\tilde X_t\}_{t\in\cT_k}$, with:
\begin{itemize}[leftmargin=1.5em,nosep]
\item a.s.\ envelope $L_X := \|\tilde X_t\|_\op\le 2R_X$ (Lemma~\ref{lem:G_bound_conf}, plus centering);
\item predictable variance
$\bigl\|\sum_{t\in\cT_k}\E[\tilde X_t^2\mid\cH_{t-1}]\bigr\|_\op\le m_k\,\sigma_*^2\,d$
(Lemma~\ref{lem:lifted_variance}, summed).
\end{itemize}
Tropp's matrix Bernstein with these inputs gives, with probability
$\ge 1-\delta$,
\[
\Bigl\|\sum_{t\in\cT_k}\tilde X_t\Bigr\|_\op
\;\le\;\sqrt{2 m_k\sigma_*^2 d\log(2d/\delta)} + \frac{2L_X\log(2d/\delta)}{3}
\;\le\;\sqrt{2 m_k\sigma_*^2 d\log(2d/\delta)} + \frac{4 R_X\log(2d/\delta)}{3},
\]
using $L_X\le 2R_X$ in the linear-envelope term. Dividing by $m_k$
and combining with the truncation-bias offset $\|\Theta_k\|_\op$ via
the same transfer step as in Theorem~\ref{thm:matrix_bernstein_conf}
yields \eqref{eq:bern_var}. The two constituent events (truncation
and Bernstein) each have mass $\ge 1-\delta$; union bound gives
$1-2\delta$.
\end{proof}

\begin{corollary}[Sharpened projector concentration; sharpening of Corollary~\ref{cor:projector_conf}]
\label{cor:projector_var}
Under the assumptions of Corollary~\ref{cor:projector_conf} together
with the additional hypotheses of Lemma~\ref{lem:lifted_variance}
($d\ge 2$, $\E[\varepsilon_t^4\mid\cH_{t-1}]\le 3\sigma_\varepsilon^4$,
indicator-truncated probes), with probability at least $1-2\delta$,
\begin{equation}
\varepsilon_k = \|\widehat P_k - P_k^\star\|_\op
\;\le\; C_{\mathrm{sub}}^\star\,\sqrt{\frac{d\,\log(2d/\delta)}{m_k}}
\;+\;\frac{16\,R_X\,\log(2d/\delta)}{3\,\lambda_{\min}\,m_k}
\;+\;\Delta_\sigma\;+\;\frac{4\|\Theta_k\|_\op}{\lambda_{\min}},
\quad
C_{\mathrm{sub}}^\star \;:=\;\frac{4\sqrt{2}\,\sigma_*}{\lambda_{\min}}.
\label{eq:proj_bound_var}
\end{equation}
The linear-in-$1/m_k$ term is dominated by the leading
$\sqrt{d/m_k}$ term whenever $m_k\gtrsim R_X^2\log(2d/\delta)/(\sigma_*^2 d)$,
which evaluates to $m_k\gtrsim d\log^2 T\cdot\log(2d/\delta) = \tilde\Theta(d)$
(with $R_X=\Theta(d\log T)$ from Lemma~\ref{lem:G_bound_conf}). The
$\tilde\Theta$ hides polylog factors of $T$ and $d/\delta$.
The sharpened constant $C_{\mathrm{sub}}^\star = \Theta(1/\lambda_{\min})$
(in the variance scale) gives leading-term scaling
$O(\sqrt{d/m_k})$, replacing the original
$C_{\mathrm{sub}}=\Theta(d)/\lambda_{\min}$ leading scaling
$O(d/\sqrt{m_k})$ — a factor-of-$\sqrt d$ improvement in the
$m_k$-dependence at any fixed $d$.
\end{corollary}
\begin{proof}
Davis--Kahan applied to \eqref{eq:bern_var}, with the leading
$\sigma_*\sqrt{2d\log(2d/\delta)/m_k}$ replacing the original
$2R_X\sqrt{\log(2d/\delta)/m_k}=\Theta(d\log T\sqrt{\log(2d/\delta)/m_k})$.
The other terms (linear $1/m_k$ Bernstein term, truncation bias,
variance-misspecification floor $\Delta_\sigma$) are unchanged.
\end{proof}

\begin{theorem}[Sharpened probe term; sharpening of Theorem~\ref{thm:spsc_regret}]
\label{thm:spsc_regret_sharp}
In addition to the assumptions of Theorem~\ref{thm:spsc_regret},
assume the conditions of Lemma~\ref{lem:lifted_variance}: $d\ge 2$
and $\E[\varepsilon_t^4\mid\cH_{t-1}]\le 3\sigma_\varepsilon^4$
(satisfied by Gaussian noise; truncation is by indicator
$\mathbf 1\{\|u_t\|\le L\}$, so $\mathcal K^{-1}$ is the full-Gaussian
inverse used throughout). Run Alg.~\ref{alg:spsc} with the sharpened
probe allocation $m_k = \lceil(B^\star\ell_k/(2A))^{2/3}\rceil$, where
\[
A := 2R_\cA S_w + c, \qquad
B^\star := 4 C_{\mathrm{sub}}^\star\,S_w\,R_\cA\,\sqrt{d\,\log(2d/\delta)}
\;=\;\Theta(\sqrt d).
\]
Assume the per-segment regime
$\ell_k\ge \ell^\star = \tilde\Theta(d)$ for all $k\in[K]$
(explicitly, $\ell^\star = C_\ell\,d\,\mathrm{polylog}(T,d/\delta)$
for a universal constant $C_\ell$, where the polylog absorbs the
$\log(2d/\delta)$ and $\log T$ factors of Corollary~\ref{cor:projector_var};
when segments are balanced, this is equivalent to
$T\ge K\ell^\star = \tilde\Theta(Kd)$).
Then, with probability $\ge 1-\delta$,
\begin{equation}
\DynReg_T^{(c)}
\;=\;\tilO\!\bigl(r\sqrt{KT}\bigr)
\;+\;\tilO\!\bigl(d^{1/3}\,K^{1/3}\,T^{2/3}\bigr)
\;+\;O(WV)\;+\;O(T\Delta_\sigma).
\label{eq:main_bound_sharp}
\end{equation}
The probe-term prefactor improves from $d^{2/3}$ to $d^{1/3}$.
\end{theorem}
\begin{proof}
The proof of Theorem~\ref{thm:spsc_regret} (\S\ref{app:proof_main})
uses $\varepsilon_k\le C_{\mathrm{sub}}\sqrt{\log(2d/\delta)/m_k}+\Delta_\sigma$
inside the step (i.c) tradeoff
\[
2R_\cA S_w\,\varepsilon_k\,n_k \;\le\; B\,n_k/\sqrt{m_k} + 2R_\cA S_w\Delta_\sigma\,n_k,
\quad B = 4 C_{\mathrm{sub}}\,S_w\,R_\cA\sqrt{\log(2d/\delta)}.
\]
Corollary~\ref{cor:projector_var} replaces this with
\[
\varepsilon_k \le C_{\mathrm{sub}}^\star\sqrt{\frac{d\log(2d/\delta)}{m_k}}
+\frac{16 R_X\log(2d/\delta)}{3\lambda_{\min}m_k}+\Delta_\sigma.
\]
Multiplying by $2R_\cA S_w n_k$, the leading variance term becomes
$B^\star n_k/\sqrt{m_k}$ with $B^\star = \Theta(\sqrt d)$ as displayed,
and the linear-Bernstein term contributes
$32 R_\cA S_w R_X \log(2d/\delta)\,n_k/(3\lambda_{\min}m_k) = \tilO(d\,n_k/m_k)$.

\emph{Threshold derivation.} From Corollary~\ref{cor:projector_var},
the linear-Bernstein contribution to the regret tradeoff is
$\tilO(d\,n_k/m_k)$ and the leading variance contribution is
$\tilO(\sqrt d\,n_k/\sqrt{m_k})$ (both up to polylog factors of
$T$ and $d/\delta$, which differ between the two terms but get
absorbed into $\tilde\Theta$ uniformly below). Their ratio at the sharpened optimum
$m_k^\star\asymp d^{1/3}\ell_k^{2/3}$ is
\[
\frac{\text{leading}}{\text{linear}}
\;\asymp\;
\frac{\sqrt{d/m_k^\star}}{d/m_k^\star}
\;=\;\sqrt{m_k^\star/d}
\;\asymp\;(\ell_k/d)^{1/3}.
\]
For domination $\ge 1$ (modulo polylog), $\ell_k\gtrsim d$ up to
polylog factors. We absorb these into $\ell^\star = \tilde\Theta(d)$;
the explicit form $C_\ell\,d\,\mathrm{polylog}(T,d/\delta)$ handles
both the $R_X = O(d\log T)$ envelope and the $\log(2d/\delta)$
Bernstein-tail factor.

\emph{Per-segment optimization.} In the regime $\ell_k\ge\ell^\star$,
the linear term is dominated, and
$\min_{m_k}\bigl(A m_k+B^\star\ell_k/\sqrt{m_k}\bigr)$
is achieved at $m_k^\star = (B^\star\ell_k/(2A))^{2/3}$ with value
\[
\Theta\bigl(A^{1/3}(B^\star)^{2/3}\bigr)\,\ell_k^{2/3} = \Theta(d^{1/3})\,\ell_k^{2/3}
\]
per segment. The leading-to-linear ratio at $m_k^\star$ is
$(B^\star\ell_k/\sqrt{m_k^\star})/(d\,n_k/m_k^\star)
\asymp (\ell_k/d)^{1/3}\to\infty$ under $\ell_k\ge\ell^\star$, confirming
domination.

\emph{Summing over segments.} Under the theorem's hypothesis
$\ell_k\ge\ell^\star$ for all $k\in[K]$, the per-segment
contributions sum via Jensen
$\sum_k\ell_k^{2/3}\le K^{1/3}T^{2/3}$ to
$\tilO(d^{1/3}K^{1/3}T^{2/3})$. (For balanced segments
$\ell_k=T/K$, the hypothesis reduces to $T\ge K\ell^\star=\tilde\Theta(Kd)$,
satisfied for $T$ large at any fixed $K,d$.)

All other steps of the proof of Theorem~\ref{thm:spsc_regret} (the
windowed-ridge in-subspace argument, the drift summation, the
variance-misspecification floor) carry over verbatim because they do
not involve $C_{\mathrm{sub}}$.\qed
\end{proof}

\begin{remark}[Algorithm structure unchanged; recommended probe rate updated]
\label{rem:no_algo_change}
Theorem~\ref{thm:spsc_regret_sharp} preserves the algorithmic
structure of Alg.~\ref{alg:spsc} (probe-then-exploit interleaving,
projector recovery via $\widehat M_k$, windowed projected ridge-UCB);
only the recommended probe rate is updated:
$m_k^\star\asymp d^{1/3}\ell_k^{2/3}$ replaces the original
$m_k^\star\asymp d^{2/3}\ell_k^{2/3}$, a $d^{1/3}$-factor reduction.
The empirical results of \S\ref{sec:experiments} were obtained with
the original schedule, so they correspond to a (theoretically)
suboptimal — but consistent — probe rate; reported regret is therefore
an upper bound on what the sharpened schedule would achieve.
The empirical $d$-scaling slope $\approx-0.16$ in Figure~\ref{fig:d-scaling}
is closer to the sharpened $d^{1/3}=+0.33$ prediction than to the
original $d^{2/3}=+0.67$, providing post-hoc empirical support for
the variance-adaptive analysis.
\end{remark}

\begin{remark}[Combined LB/UB picture, after both extensions]
\label{rem:combined_after_ext}
Theorem~\ref{thm:spsc_lb_piecewise} and
Theorem~\ref{thm:spsc_regret_sharp} together yield
\[
\Omega\!\bigl(c^{1/3}\,K^{1/3}T^{2/3}\bigr)
\;\le\;\text{(probe-acquisition regret of SPSC)}\;\le\;
\tilO\!\bigl(d^{1/3}\,K^{1/3}T^{2/3}\bigr),
\]
tight in $K$ and $T$, with a residual $d^{1/3}$ gap that the
empirical $d$-scaling fit (Figure~\ref{fig:d-scaling}) suggests is
also loose; closing it likely requires the subspace-restricted
concentration sketched in Remark~\ref{rem:dim_dep}, which is beyond
the scope of this purely-additive sharpening.
\end{remark}

% =====================================================================
% End extensions.tex
% =====================================================================
` to conf.tex just before `\end{document}`
%     (i.e., after Section J / Table tab:published_comparison).
%   * Optionally add an entry to the appendix TOC (line ~847 of conf.tex):
%       \tocentry{K}{app:extensions}{Extensions: piecewise LB and sharpened d-dependence}
% =====================================================================

\section{Extensions: piecewise lower bound and sharpened \texorpdfstring{$d$}{d}-dependence}
\label{app:extensions}

This section supplements (without modifying) Theorems~\ref{thm:spsc_regret}--\ref{thm:spsc_lb}
with two additive results: a piecewise-stationary lower bound that
matches the upper-bound rate in $K$ (\S\ref{app:ext_piecewise_lb}),
and a sharpened $d$-dependence in the probe term via variance-adaptive
matrix Bernstein (\S\ref{app:ext_sharpened_d}). The original theorems
remain in force; the results below are stated as \emph{independent}
sharpenings.

\subsection{Piecewise-stationary lower bound}
\label{app:ext_piecewise_lb}

Theorem~\ref{thm:spsc_lb} establishes the LB only for $K=1$ (single
segment of length $T$). The $K$-segment extension follows by applying
the single-segment bound conditionally inside each segment, exploiting
the fact that the per-segment hidden direction $v_k$ in the
construction is independent of all previous segments. No new technical
tools beyond Theorem~\ref{thm:spsc_lb} itself are required.

\begin{theorem}[Piecewise lower bound; extension of Theorem~\ref{thm:spsc_lb} to $K\ge 1$]
\label{thm:spsc_lb_piecewise}
Let $\Pi$ be the policy class of Theorem~\ref{thm:spsc_lb} (bounded
exploit actions $\|x_t\|\le R_\cA$, isotropic Gaussian probes
$u_t\sim\cN(0,I_d)$). Fix $K\ge 1$ and segment lengths
$\{\ell_k\}_{k=1}^K$ with $\sum_k\ell_k=T$. There exist constants
$c_0,d_0,\ell_0>0$ depending only on $R_\cA,\sigma_\varepsilon,c$
(in fact, the same constants as Theorem~\ref{thm:spsc_lb}) such that
whenever $\min_k\ell_k\ge\ell_0$ and $d\ge d_0\,(\min_k\ell_k)^{1/3}$,
\begin{equation}
\inf_{\pi\in\Pi}\sup_\nu\ \E^\pi_\nu[\DynReg_T^{(c)}]
\;\ge\; c_0\,c^{1/3}(R_\cA\sigma_\varepsilon)^{2/3}\sum_{k=1}^K\ell_k^{2/3}.
\label{eq:lb_piecewise}
\end{equation}
For uniform segments $\ell_k=T/K$ this yields
\begin{equation}
\inf_{\pi\in\Pi}\sup_\nu\ \E^\pi_\nu[\DynReg_T^{(c)}]
\;\ge\; c_0\,c^{1/3}(R_\cA\sigma_\varepsilon)^{2/3}\,K^{1/3}T^{2/3},
\label{eq:lb_piecewise_uniform}
\end{equation}
matching the probe-acquisition rate of Theorem~\ref{thm:spsc_regret}
in both $K$ and $T$ (with the same residual $d$-prefactor gap discussed
in Remark~\ref{rem:dim_dep}, sharpened in
\S\ref{app:ext_sharpened_d} below).
\end{theorem}

\begin{proof}
\emph{Construction.} Let $v_1,\dots,v_K$ be i.i.d.\ uniform on
$\mathbb S^{d-1}$, drawn independently of the policy. For
per-segment scales $\alpha_k>0$ to be specified, define the instance
$\nu(\vec v):\theta_t = \alpha_k v_k$ for $t\in\cI_k$. Let
$\nu_0:\theta_t\equiv 0$ denote the global null. The per-segment
``zero-on-segment-$k$'' alternative is
$\nu^{(-k)}(\vec v):\theta_t = \alpha_{k'}v_{k'}\mathbf 1\{k'\ne k\}$
for $t\in\cI_{k'}$ — i.e., identical to $\nu(\vec v)$ except on
segment $k$ where $\theta_t=0$.

\emph{Per-segment KL.} The trajectory laws under $\nu(\vec v)$ and
$\nu^{(-k)}(\vec v)$ agree on segments $k'\ne k$ and differ only on
segment $k$. The sequential KL chain rule therefore gives
\[
\mathrm{KL}\bigl(\mathbb P_{\nu^{(-k)}(\vec v)}\,\bigl\|\,\mathbb P_{\nu(\vec v)}\bigr)
\;=\; \frac{\alpha_k^2}{2\sigma_\varepsilon^2}\,
\E^\pi_{\nu^{(-k)}(\vec v)}\!\Bigl[\sum_{t\in\cI_k}(x_t^\top v_k)^2\Bigr].
\]
Under $\nu^{(-k)}(\vec v)$, on segment $k$ the parameter is $0$, so
the segment-$k$ marginal trajectory is independent of $v_k$.
Consequently $v_k$ is independent of the actions $\{x_t\}_{t\in\cI_k}$
under $\nu^{(-k)}(\vec v)$, and the trace-identity argument of
Lemma~\ref{lem:trace_kl_conf} applies verbatim with horizon $\ell_k$,
probe count $m_k$, and exploit count $n_k=\ell_k-m_k$:
\[
\E_{v_k}\E^\pi_{\nu^{(-k)}(\vec v)}\!\Bigl[\sum_{t\in\cI_k}(x_t^\top v_k)^2\Bigr]
\;\le\; m_k + R_\cA^2 n_k/d \;=:\;Q_k,
\]
where the average is over $v_k\sim\mathrm{Unif}(\mathbb S^{d-1})$. As in
Lemma~\ref{lem:trace_kl_conf}, the existence of an extremal direction
gives a deterministic $v_k^\star\in\mathbb S^{d-1}$ for which the
inequality holds with $v_k^\star$ replacing the average. Setting
$\alpha_k^2 := \sigma_\varepsilon^2/(4 Q_k)$ yields
\[
\mathrm{KL}\bigl(\mathbb P_{\nu^{(-k)}(\vec v^\star)}\,\bigl\|\,\mathbb P_{\nu(\vec v^\star)}\bigr)
\;\le\;\frac{\alpha_k^2 Q_k}{2\sigma_\varepsilon^2}\;=\;\frac{1}{8},
\]
hence Pinsker gives per-segment TV $\le 1/4$.

\emph{Per-segment regret transfer.} The optimal action under
$\nu(\vec v^\star)$ on segment $k$ is $x_t^\star=R_\cA v_k^\star$ with
reward $\alpha_k R_\cA$; under $\nu^{(-k)}(\vec v^\star)$ the
parameter on segment $k$ is $0$ so any action is optimal, and the
segment-$k$ regret of the policy is identically $0$. The
arm-regret-vs.-alignment manipulation of \S\ref{app:proof_lb}
(equation~\eqref{eq:reg_to_alignment_conf}), applied to the segment-$k$
exploit rounds, gives
\[
\E^\pi_{\nu(\vec v^\star)}\!\Bigl[\sum_{t\in\cI_k\cap\cT_{\mathrm{ex}}}\Delta_t\Bigr]
\;\ge\;\alpha_k n_k R_\cA\;-\;\alpha_k\,\E^\pi_{\nu(\vec v^\star)}[f_k],\quad
f_k:=\sum_{t\in\cI_k\cap\cT_{\mathrm{ex}}}|v_k^{\star\top}x_t|.
\]
TV transfer from $\nu(\vec v^\star)$ to $\nu^{(-k)}(\vec v^\star)$
with $\|f_k\|_\infty\le n_k R_\cA$ and segment-$k$
$d_{\mathrm{TV}}\le 1/4$:
\[
\E^\pi_{\nu(\vec v^\star)}[f_k]\;\le\;\E^\pi_{\nu^{(-k)}(\vec v^\star)}[f_k]+\tfrac12\,n_k R_\cA.
\]
Under $\nu^{(-k)}(\vec v^\star)$, by Cauchy--Schwarz and the
trace-identity bound applied on segment $k$,
$\E^\pi_{\nu^{(-k)}(\vec v^\star)}[f_k]\le R_\cA n_k/\sqrt d$.
Combining (with $d\ge 16$ so $1/\sqrt d\le 1/4$):
\begin{equation}
\E^\pi_{\nu(\vec v^\star)}\!\Bigl[\sum_{t\in\cI_k\cap\cT_{\mathrm{ex}}}\Delta_t\Bigr]
\;\ge\;\tfrac14\,\alpha_k n_k R_\cA
\;=\;\frac{R_\cA\sigma_\varepsilon\,n_k}{8\sqrt{Q_k}}.
\label{eq:per_seg_regret_lb}
\end{equation}

\emph{Per-segment costed regret.} Adding the segment-$k$ probe cost
$cm_k$,
\[
\E^\pi_{\nu(\vec v^\star)}[\DynReg^{(c)}_k]
\;\ge\; cm_k + \frac{R_\cA\sigma_\varepsilon\,n_k}{8\sqrt{Q_k}},
\]
with the same functional form as the single-segment bound
\eqref{eq:dynreg_lower_conf}. The Regime-I/Regime-II optimization of
\S\ref{app:proof_lb} applies verbatim with $T\to\ell_k$, yielding
\[
\E^\pi_{\nu(\vec v^\star)}[\DynReg^{(c)}_k]
\;\ge\;c_0\,c^{1/3}(R_\cA\sigma_\varepsilon)^{2/3}\ell_k^{2/3}
\]
in the regime $\ell_k\ge\ell_0$ and $d\ge d_0\,\ell_k^{1/3}$, with
$c_0,d_0,\ell_0$ as in Theorem~\ref{thm:spsc_lb}.

\emph{Summing across segments.} Total dynamic regret is the sum of
per-segment costed regrets (probes and exploit rounds partition by
segment):
\[
\E^\pi_{\nu(\vec v^\star)}[\DynReg^{(c)}_T]
\;=\;\sum_{k=1}^K\E^\pi_{\nu(\vec v^\star)}[\DynReg^{(c)}_k]
\;\ge\;c_0\,c^{1/3}(R_\cA\sigma_\varepsilon)^{2/3}\sum_{k=1}^K\ell_k^{2/3}.
\]
Taking $\sup_\nu\ge$ this specific instance yields
\eqref{eq:lb_piecewise}; specialising to $\ell_k=T/K$ gives
$\sum_k\ell_k^{2/3}=K\,(T/K)^{2/3}=K^{1/3}T^{2/3}$, hence
\eqref{eq:lb_piecewise_uniform}.\qed
\end{proof}

\begin{remark}[The argument uses only Theorem~\ref{thm:spsc_lb}'s tools]
\label{rem:lb_piecewise_tools}
The proof reuses the trace-identity Lemma~\ref{lem:trace_kl_conf} and
the Le~Cam transfer of \S\ref{app:proof_lb} per segment. The only new
ingredient is the observation that the segment-$k$ KL between
$\nu(\vec v^\star)$ and $\nu^{(-k)}(\vec v^\star)$ is exactly the
single-segment KL with parameters $(\ell_k, m_k, n_k)$, because the
two laws agree outside segment $k$. No multi-hypothesis Fano, no
sliding-window reduction, and no new concentration is invoked.
\end{remark}

\begin{remark}[Adaptive probe schedules]
\label{rem:lb_piecewise_adaptive}
The bound \eqref{eq:per_seg_regret_lb} is in terms of the policy's
\emph{realised} segment-$k$ probe count $m_k$ and exploit count $n_k$,
which may be adaptive (chosen based on history through $\tau_{k-1}$).
The per-segment Regime-I/Regime-II optimization minimizes
$cm_k+R_\cA\sigma_\varepsilon n_k/(8\sqrt{Q_k})$ over $m_k\in[0,\ell_k]$;
the policy's adaptive choice of $m_k$ can only weakly increase this
quantity. Hence the LB applies under any (possibly adaptive)
per-segment probe allocation.
\end{remark}

\begin{remark}[Combined LB/UB picture]
\label{rem:combined_lb_ub}
Theorem~\ref{thm:spsc_lb_piecewise} combined with
Theorem~\ref{thm:spsc_regret} yields the matching pair
\[
\Omega(c^{1/3}\,K^{1/3}T^{2/3})
\;\le\;\text{(probe-acquisition regret)}\;\le\;
\tilO(d^{2/3}\,K^{1/3}T^{2/3}),
\]
tight in $K$ and $T$. The remaining gap is the $d$-prefactor on the
upper-bound side; \S\ref{app:ext_sharpened_d} narrows that gap from
$d^{2/3}$ to $d^{1/3}$ via a sharper concentration argument, leaving
a residual $d^{1/3}$ open question.
\end{remark}

\subsection{Sharpened \texorpdfstring{$d$}{d}-dependence via variance-adaptive Bernstein}
\label{app:ext_sharpened_d}

Theorem~\ref{thm:matrix_bernstein_conf} (matrix Freedman) bounds
$\|\widehat M_k-\bar M_k^{\mathrm{probe}}-\widetilde B\|_\op$ using the
worst-case envelope $\|\tilde X_t\|_\op\le 2R_X$ with
$R_X=\Theta(L^2)=\Theta(d\log T)$. The leading term
$2R_X\sqrt{\log(2d/\delta)/m_k}$ then yields
$C_{\mathrm{sub}}=8R_X/\lambda_{\min}=\Theta(d\log T/\lambda_{\min})$
in Corollary~\ref{cor:projector_conf}, and the constant
$B = 4 C_{\mathrm{sub}} S_w R_\cA\sqrt{\log(2d/\delta)} = \Theta(d)$
propagates through the proof of Theorem~\ref{thm:spsc_regret} to give
the $\tilO(d^{2/3}K^{1/3}T^{2/3})$ probe-acquisition rate.

The envelope is loose. For isotropic Gaussian $u_t$,
$\|u_t\|^2$ concentrates near $d$ with $\sqrt d$ fluctuations, so the
\emph{predictable second moment} of $\tilde G_t$ in operator norm is
$\Theta(d)$ per round, not $\Theta(d^2)$ as the envelope suggests.
The variance-adaptive matrix Bernstein
inequality~\citep[Theorem~7.7.1]{tropp2015introduction} exploits this
gap and yields a sharper concentration. The algorithmic structure of
SPSC (probe-then-exploit interleaving, projector recovery, projected
ridge-UCB) is unchanged; only the recommended probe rate $m_k$ is
updated by a $d^{1/3}$ factor
(from $m_k^\star\asymp d^{2/3}\ell_k^{2/3}$ to
$m_k^\star\asymp d^{1/3}\ell_k^{2/3}$; see
Theorem~\ref{thm:spsc_regret_sharp} and Remark~\ref{rem:no_algo_change}).

\begin{lemma}[Predictable second moment of the lifted probe sample]
\label{lem:lifted_variance}
Assume $d\ge 2$, $\|\theta_t\|\le S_w$ a.s., and that the noise has a
bounded fourth moment $\E[\varepsilon_t^4\mid\cH_{t-1}]\le 3\sigma_\varepsilon^4$
(satisfied by Gaussian $\varepsilon_t\sim\cN(0,\sigma_\varepsilon^2)$).
Let $u_t\sim\cN(0,I_d)$ on probe rounds (truncation by an indicator
$\mathbf 1\{\|u_t\|\le L\}$ only decreases the second moment in the
PSD order, so the bound below applies; if probes are instead
distributed $u_t\mid\{\|u_t\|\le L\}$, the same bound holds up to a
$1+o(1)$ factor as $L\to\infty$). Define $G_t := \mathcal K^{-1}(s_t u_t u_t^\top)$
and $\tilde X_t := G_t\mathbf 1\{\cA_t\}-\E[G_t\mathbf 1\{\cA_t\}\mid\cH_{t-1}]$.
Then there exists $\sigma_*^2$ depending only on $S_w,\sigma_\varepsilon,\hat\sigma$
(and not on $d$) such that
\begin{equation}
\bigl\|\E[\tilde X_t^2\mid\cH_{t-1}]\bigr\|_\op
\;\le\;\sigma_*^2\,d.
\label{eq:lifted_var}
\end{equation}
Explicitly, $\sigma_*^2 \le C\,(S_w^4 + \sigma_\varepsilon^4 + \hat\sigma^4)$
with $C=100$ (using the exact trace-expansion bound established in the
proof; the value of $C$ does not affect the leading-order rate).
\end{lemma}
\begin{proof}
Centering only reduces the second moment in the PSD order:
$\E[\tilde X_t^2\mid\cH_{t-1}]\preceq\E[G_t^2\mathbf 1\{\cA_t\}\mid\cH_{t-1}]\preceq\E[G_t^2\mid\cH_{t-1}]$.
We bound $\|\E[G_t^2\mid\cH_{t-1}]\|_\op$ by conditioning further on
$\theta_t$ and proving a bound uniform in $\theta_t$ with
$\|\theta_t\|\le S_w$. By tower-property the same bound then applies
to $\E[\,\cdot\mid\cH_{t-1}]$.

By Lemma~\ref{lem:K_inverse},
\(
G_t = \tfrac{s_t}{2}\,u_t u_t^\top - \tfrac{s_t\|u_t\|^2}{2(d{+}2)}\,I.
\)
Squaring and using $(uu^\top)^2 = \|u\|^2 uu^\top$:
\begin{equation}
G_t^2 = \frac{d}{4(d+2)}\,s_t^2\|u_t\|^2\, u_t u_t^\top
       \;+\;\frac{s_t^2\,\|u_t\|^4}{4(d+2)^2}\,I.
\label{eq:Gt_squared}
\end{equation}

\emph{Leading-in-$d$ block.}
Conditional on $(\cH_{t-1},\theta_t)$, define $e_1 := \theta_t/\|\theta_t\|$
(any unit vector if $\theta_t=0$). This $e_1$ is now deterministic, so
the decomposition $u_t = z_t e_1 + u_{t,\perp}$ is well-defined with
$z_t := u_t^\top e_1\sim\cN(0,1)$ and $u_{t,\perp}$ independent of
$z_t$, distributed as $\cN(0,I-e_1e_1^\top)$ on the $(d{-}1)$-dim
subspace $e_1^\perp$. Crucially, $s_t = (z_t\|\theta_t\|+\varepsilon_t)^2-\hat\sigma^2$
depends on $u_t$ only through $z_t$, so $s_t$ is independent of
$u_{t,\perp}$ given $(\cH_{t-1},\theta_t)$.

Expand $u_tu_t^\top = z_t^2 e_1e_1^\top + z_t(e_1 u_{t,\perp}^\top + u_{t,\perp}e_1^\top) + u_{t,\perp}u_{t,\perp}^\top$
and $\|u_t\|^2 = z_t^2 + \|u_{t,\perp}\|^2$. The cross terms $z_t e_1 u_{t,\perp}^\top$
vanish under $\E[\cdot\mid\cH_{t-1},\theta_t]$ by symmetry of
$u_{t,\perp}$. The surviving contributions split into:
\begin{align*}
\E\bigl[s_t^2\|u_t\|^2\,z_t^2\,e_1e_1^\top \,\big|\,\cH_{t-1},\theta_t\bigr]
&= \Bigl(\E[s_t^2 z_t^4] + (d{-}1)\,\E[s_t^2 z_t^2]\Bigr)\,e_1e_1^\top, \\
\E\bigl[s_t^2\|u_t\|^2\,u_{t,\perp}u_{t,\perp}^\top \,\big|\,\cH_{t-1},\theta_t\bigr]
&= \E[s_t^2 z_t^2]\,(I-e_1e_1^\top) + \E[s_t^2]\,(d{+}1)(I-e_1e_1^\top),
\end{align*}
where conditional expectations are abbreviated, and we used:
\begin{itemize}[leftmargin=1.5em,nosep]
\item $\E[u_{t,\perp}u_{t,\perp}^\top]=I-e_1e_1^\top$ (covariance);
\item $\E[\|u_{t,\perp}\|^2]=d-1$ ($(d{-}1)$-dim chi-square);
\item $\E[\|u_{t,\perp}\|^2\,u_{t,\perp}u_{t,\perp}^\top]=(d+1)(I-e_1e_1^\top)$
  (Isserlis on the $(d{-}1)$-dim subspace, valid for $d\ge 2$);
\item independence of $u_{t,\perp}$ from $s_t$ given $(\cH_{t-1},\theta_t)$
  splits the joint expectation in the second display.
\end{itemize}

\emph{Bounding the conditional moments of $s_t$.} Write
$\xi_t := z_t\|\theta_t\|+\varepsilon_t$, so $s_t = \xi_t^2-\hat\sigma^2$ and
$s_t^2 = \xi_t^4 - 2\hat\sigma^2\xi_t^2 + \hat\sigma^4$. Given
$\theta_t$, $\xi_t$ has variance $\|\theta_t\|^2+\sigma_\varepsilon^2\le S_w^2+\sigma_\varepsilon^2$.
Using the assumed fourth-moment bound $\E[\varepsilon_t^4\mid\cH_{t-1}]\le 3\sigma_\varepsilon^4$
and Gaussian moments $\E[z_t^k]=k!!$ for $k\le 8$ (in particular
$\E[z_t^8]=105$ — needed because $s_t^2 z_t^4$ contains a $z_t^8$
contribution from $\xi_t^4 z_t^4 = (z_t\|\theta_t\|+\varepsilon_t)^4 z_t^4$),
straightforward expansion gives
\[
\E[s_t^2\mid\cH_{t-1},\theta_t]\le C_0\,D,\quad
\E[s_t^2 z_t^2\mid\cH_{t-1},\theta_t]\le C_0\,D,\quad
\E[s_t^2 z_t^4\mid\cH_{t-1},\theta_t]\le C_0\,D,
\]
with $D := S_w^4+\sigma_\varepsilon^4+\hat\sigma^4$ and a universal
constant $C_0$ (the dominant contribution is
$\E[\xi_t^4 z_t^4]\le 105\,S_w^4 + 90\,S_w^2\sigma_\varepsilon^2+9\sigma_\varepsilon^4
\le 320\,D$ via $\E[z_t^8]=105$ and AM-GM, so $C_0\le 320$ suffices).

\emph{Combining the blocks.} The parallel-block coefficient is
$A_1 := \E[s_t^2 z_t^4]+(d-1)\E[s_t^2 z_t^2]\le d\,C_0\,D$, and the
perpendicular-block coefficient is
$A_2 := \E[s_t^2 z_t^2]+(d+1)\E[s_t^2]\le (d+2)\,C_0\,D$. Hence
$\|\E[s_t^2\|u_t\|^2 u_tu_t^\top\mid\cH_{t-1},\theta_t]\|_\op
\le \max(A_1,A_2)\le (d+2)\,C_0\,D$. Multiplying by the
$G_t^2$-prefactor $d/(4(d+2))$ from \eqref{eq:Gt_squared} yields a
leading-block contribution
$\le (d/(4(d+2)))\cdot (d+2)\,C_0\,D = (C_0/4)\,d\,D$ exactly.

\emph{Trace correction (exact bound).} The second term of
\eqref{eq:Gt_squared} contributes
$\E[s_t^2\|u_t\|^4]/(4(d+2)^2)\cdot I$. Expanding
$\|u_t\|^4 = z_t^4 + 2z_t^2\|u_{t,\perp}\|^2 + \|u_{t,\perp}\|^4$ and
using $\E[\|u_{t,\perp}\|^4]=(d-1)(d+1)$,
\[
\E[s_t^2\|u_t\|^4\mid\cH_{t-1},\theta_t]
\;\le\;C_0\,D\,\bigl(1+2(d-1)+(d^2-1)\bigr)
\;=\;C_0\,D\,(d^2+2d-2).
\]
Divided by $4(d+2)^2$, the trace contribution to
$\|\E[G_t^2\mid\cH_{t-1},\theta_t]\|_\op$ is at most
$C_0\,D\,(d^2+2d-2)/(4(d+2)^2)$.

\emph{Conclusion (exact prefactor accounting).} Combining the
leading-block contribution $\le (C_0/4)\,d\,D$ with the trace
correction yields, after dividing by $d\,D$,
\[
\frac{\|\E[G_t^2\mid\cH_{t-1},\theta_t]\|_\op}{d\,D}
\;\le\;\frac{C_0}{4}
\;+\;\frac{C_0\,(d^2+2d-2)}{4d\,(d+2)^2}.
\]
With $C_0=320$, the second term is monotone decreasing in $d$ for
$d\ge 2$: at $d=2$, $C_0\cdot 6/(8\cdot 16)=15$; at $d=4$,
$C_0\cdot 22/(16\cdot 36)\approx 12.22$; as $d\to\infty$, the term
approaches $0$. Combined with the constant first term $C_0/4=80$,
the per-$dD$ bound is $\le 80+15=95$ at $d=2$, $\le 92.22$ at $d=4$,
and $\to 80$ as $d\to\infty$ — uniformly $\le 100$. Hence
$\sigma_*^2 = C\,D$ with $C = 100$ suffices for all $d\ge 2$.
Taking conditional expectation over $\theta_t$ preserves the bound
by linearity and operator-norm convexity, yielding
\eqref{eq:lifted_var}.
\end{proof}

\begin{theorem}[Variance-adaptive concentration; sharpening of Theorem~\ref{thm:matrix_bernstein_conf}]
\label{thm:matrix_bernstein_var}
Under the assumptions of Theorem~\ref{thm:matrix_bernstein_conf}
together with the additional hypotheses of Lemma~\ref{lem:lifted_variance}
($d\ge 2$ and $\E[\varepsilon_t^4\mid\cH_{t-1}]\le 3\sigma_\varepsilon^4$;
truncation by indicator $\mathbf 1\{\|u_t\|\le L\}$), with
probability at least $1-2\delta$ (jointly over the truncation event of
Lemma~\ref{lem:G_bound_conf} and the Bernstein event),
\begin{equation}
\bigl\|\widehat M_k-\bar M_k^{\mathrm{probe}}-\widetilde B\bigr\|_\op
\;\le\;\sigma_*\sqrt{\frac{2\,d\,\log(2d/\delta)}{m_k}}
\;+\;\frac{4 R_X\log(2d/\delta)}{3\,m_k}
\;+\;\|\Theta_k\|_\op,
\label{eq:bern_var}
\end{equation}
where $R_X = \Theta(L^2)=\Theta(d\log T)$ is the worst-case envelope of
Lemma~\ref{lem:G_bound_conf} and $\sigma_*^2$ is the variance constant
of Lemma~\ref{lem:lifted_variance}.
\end{theorem}
\begin{proof}
This is matrix Bernstein \citep[Theorem~7.7.1]{tropp2015introduction}
applied to the centered MDS $\{\tilde X_t\}_{t\in\cT_k}$, with:
\begin{itemize}[leftmargin=1.5em,nosep]
\item a.s.\ envelope $L_X := \|\tilde X_t\|_\op\le 2R_X$ (Lemma~\ref{lem:G_bound_conf}, plus centering);
\item predictable variance
$\bigl\|\sum_{t\in\cT_k}\E[\tilde X_t^2\mid\cH_{t-1}]\bigr\|_\op\le m_k\,\sigma_*^2\,d$
(Lemma~\ref{lem:lifted_variance}, summed).
\end{itemize}
Tropp's matrix Bernstein with these inputs gives, with probability
$\ge 1-\delta$,
\[
\Bigl\|\sum_{t\in\cT_k}\tilde X_t\Bigr\|_\op
\;\le\;\sqrt{2 m_k\sigma_*^2 d\log(2d/\delta)} + \frac{2L_X\log(2d/\delta)}{3}
\;\le\;\sqrt{2 m_k\sigma_*^2 d\log(2d/\delta)} + \frac{4 R_X\log(2d/\delta)}{3},
\]
using $L_X\le 2R_X$ in the linear-envelope term. Dividing by $m_k$
and combining with the truncation-bias offset $\|\Theta_k\|_\op$ via
the same transfer step as in Theorem~\ref{thm:matrix_bernstein_conf}
yields \eqref{eq:bern_var}. The two constituent events (truncation
and Bernstein) each have mass $\ge 1-\delta$; union bound gives
$1-2\delta$.
\end{proof}

\begin{corollary}[Sharpened projector concentration; sharpening of Corollary~\ref{cor:projector_conf}]
\label{cor:projector_var}
Under the assumptions of Corollary~\ref{cor:projector_conf} together
with the additional hypotheses of Lemma~\ref{lem:lifted_variance}
($d\ge 2$, $\E[\varepsilon_t^4\mid\cH_{t-1}]\le 3\sigma_\varepsilon^4$,
indicator-truncated probes), with probability at least $1-2\delta$,
\begin{equation}
\varepsilon_k = \|\widehat P_k - P_k^\star\|_\op
\;\le\; C_{\mathrm{sub}}^\star\,\sqrt{\frac{d\,\log(2d/\delta)}{m_k}}
\;+\;\frac{16\,R_X\,\log(2d/\delta)}{3\,\lambda_{\min}\,m_k}
\;+\;\Delta_\sigma\;+\;\frac{4\|\Theta_k\|_\op}{\lambda_{\min}},
\quad
C_{\mathrm{sub}}^\star \;:=\;\frac{4\sqrt{2}\,\sigma_*}{\lambda_{\min}}.
\label{eq:proj_bound_var}
\end{equation}
The linear-in-$1/m_k$ term is dominated by the leading
$\sqrt{d/m_k}$ term whenever $m_k\gtrsim R_X^2\log(2d/\delta)/(\sigma_*^2 d)$,
which evaluates to $m_k\gtrsim d\log^2 T\cdot\log(2d/\delta) = \tilde\Theta(d)$
(with $R_X=\Theta(d\log T)$ from Lemma~\ref{lem:G_bound_conf}). The
$\tilde\Theta$ hides polylog factors of $T$ and $d/\delta$.
The sharpened constant $C_{\mathrm{sub}}^\star = \Theta(1/\lambda_{\min})$
(in the variance scale) gives leading-term scaling
$O(\sqrt{d/m_k})$, replacing the original
$C_{\mathrm{sub}}=\Theta(d)/\lambda_{\min}$ leading scaling
$O(d/\sqrt{m_k})$ — a factor-of-$\sqrt d$ improvement in the
$m_k$-dependence at any fixed $d$.
\end{corollary}
\begin{proof}
Davis--Kahan applied to \eqref{eq:bern_var}, with the leading
$\sigma_*\sqrt{2d\log(2d/\delta)/m_k}$ replacing the original
$2R_X\sqrt{\log(2d/\delta)/m_k}=\Theta(d\log T\sqrt{\log(2d/\delta)/m_k})$.
The other terms (linear $1/m_k$ Bernstein term, truncation bias,
variance-misspecification floor $\Delta_\sigma$) are unchanged.
\end{proof}

\begin{theorem}[Sharpened probe term; sharpening of Theorem~\ref{thm:spsc_regret}]
\label{thm:spsc_regret_sharp}
In addition to the assumptions of Theorem~\ref{thm:spsc_regret},
assume the conditions of Lemma~\ref{lem:lifted_variance}: $d\ge 2$
and $\E[\varepsilon_t^4\mid\cH_{t-1}]\le 3\sigma_\varepsilon^4$
(satisfied by Gaussian noise; truncation is by indicator
$\mathbf 1\{\|u_t\|\le L\}$, so $\mathcal K^{-1}$ is the full-Gaussian
inverse used throughout). Run Alg.~\ref{alg:spsc} with the sharpened
probe allocation $m_k = \lceil(B^\star\ell_k/(2A))^{2/3}\rceil$, where
\[
A := 2R_\cA S_w + c, \qquad
B^\star := 4 C_{\mathrm{sub}}^\star\,S_w\,R_\cA\,\sqrt{d\,\log(2d/\delta)}
\;=\;\Theta(\sqrt d).
\]
Assume the per-segment regime
$\ell_k\ge \ell^\star = \tilde\Theta(d)$ for all $k\in[K]$
(explicitly, $\ell^\star = C_\ell\,d\,\mathrm{polylog}(T,d/\delta)$
for a universal constant $C_\ell$, where the polylog absorbs the
$\log(2d/\delta)$ and $\log T$ factors of Corollary~\ref{cor:projector_var};
when segments are balanced, this is equivalent to
$T\ge K\ell^\star = \tilde\Theta(Kd)$).
Then, with probability $\ge 1-\delta$,
\begin{equation}
\DynReg_T^{(c)}
\;=\;\tilO\!\bigl(r\sqrt{KT}\bigr)
\;+\;\tilO\!\bigl(d^{1/3}\,K^{1/3}\,T^{2/3}\bigr)
\;+\;O(WV)\;+\;O(T\Delta_\sigma).
\label{eq:main_bound_sharp}
\end{equation}
The probe-term prefactor improves from $d^{2/3}$ to $d^{1/3}$.
\end{theorem}
\begin{proof}
The proof of Theorem~\ref{thm:spsc_regret} (\S\ref{app:proof_main})
uses $\varepsilon_k\le C_{\mathrm{sub}}\sqrt{\log(2d/\delta)/m_k}+\Delta_\sigma$
inside the step (i.c) tradeoff
\[
2R_\cA S_w\,\varepsilon_k\,n_k \;\le\; B\,n_k/\sqrt{m_k} + 2R_\cA S_w\Delta_\sigma\,n_k,
\quad B = 4 C_{\mathrm{sub}}\,S_w\,R_\cA\sqrt{\log(2d/\delta)}.
\]
Corollary~\ref{cor:projector_var} replaces this with
\[
\varepsilon_k \le C_{\mathrm{sub}}^\star\sqrt{\frac{d\log(2d/\delta)}{m_k}}
+\frac{16 R_X\log(2d/\delta)}{3\lambda_{\min}m_k}+\Delta_\sigma.
\]
Multiplying by $2R_\cA S_w n_k$, the leading variance term becomes
$B^\star n_k/\sqrt{m_k}$ with $B^\star = \Theta(\sqrt d)$ as displayed,
and the linear-Bernstein term contributes
$32 R_\cA S_w R_X \log(2d/\delta)\,n_k/(3\lambda_{\min}m_k) = \tilO(d\,n_k/m_k)$.

\emph{Threshold derivation.} From Corollary~\ref{cor:projector_var},
the linear-Bernstein contribution to the regret tradeoff is
$\tilO(d\,n_k/m_k)$ and the leading variance contribution is
$\tilO(\sqrt d\,n_k/\sqrt{m_k})$ (both up to polylog factors of
$T$ and $d/\delta$, which differ between the two terms but get
absorbed into $\tilde\Theta$ uniformly below). Their ratio at the sharpened optimum
$m_k^\star\asymp d^{1/3}\ell_k^{2/3}$ is
\[
\frac{\text{leading}}{\text{linear}}
\;\asymp\;
\frac{\sqrt{d/m_k^\star}}{d/m_k^\star}
\;=\;\sqrt{m_k^\star/d}
\;\asymp\;(\ell_k/d)^{1/3}.
\]
For domination $\ge 1$ (modulo polylog), $\ell_k\gtrsim d$ up to
polylog factors. We absorb these into $\ell^\star = \tilde\Theta(d)$;
the explicit form $C_\ell\,d\,\mathrm{polylog}(T,d/\delta)$ handles
both the $R_X = O(d\log T)$ envelope and the $\log(2d/\delta)$
Bernstein-tail factor.

\emph{Per-segment optimization.} In the regime $\ell_k\ge\ell^\star$,
the linear term is dominated, and
$\min_{m_k}\bigl(A m_k+B^\star\ell_k/\sqrt{m_k}\bigr)$
is achieved at $m_k^\star = (B^\star\ell_k/(2A))^{2/3}$ with value
\[
\Theta\bigl(A^{1/3}(B^\star)^{2/3}\bigr)\,\ell_k^{2/3} = \Theta(d^{1/3})\,\ell_k^{2/3}
\]
per segment. The leading-to-linear ratio at $m_k^\star$ is
$(B^\star\ell_k/\sqrt{m_k^\star})/(d\,n_k/m_k^\star)
\asymp (\ell_k/d)^{1/3}\to\infty$ under $\ell_k\ge\ell^\star$, confirming
domination.

\emph{Summing over segments.} Under the theorem's hypothesis
$\ell_k\ge\ell^\star$ for all $k\in[K]$, the per-segment
contributions sum via Jensen
$\sum_k\ell_k^{2/3}\le K^{1/3}T^{2/3}$ to
$\tilO(d^{1/3}K^{1/3}T^{2/3})$. (For balanced segments
$\ell_k=T/K$, the hypothesis reduces to $T\ge K\ell^\star=\tilde\Theta(Kd)$,
satisfied for $T$ large at any fixed $K,d$.)

All other steps of the proof of Theorem~\ref{thm:spsc_regret} (the
windowed-ridge in-subspace argument, the drift summation, the
variance-misspecification floor) carry over verbatim because they do
not involve $C_{\mathrm{sub}}$.\qed
\end{proof}

\begin{remark}[Algorithm structure unchanged; recommended probe rate updated]
\label{rem:no_algo_change}
Theorem~\ref{thm:spsc_regret_sharp} preserves the algorithmic
structure of Alg.~\ref{alg:spsc} (probe-then-exploit interleaving,
projector recovery via $\widehat M_k$, windowed projected ridge-UCB);
only the recommended probe rate is updated:
$m_k^\star\asymp d^{1/3}\ell_k^{2/3}$ replaces the original
$m_k^\star\asymp d^{2/3}\ell_k^{2/3}$, a $d^{1/3}$-factor reduction.
The empirical results of \S\ref{sec:experiments} were obtained with
the original schedule, so they correspond to a (theoretically)
suboptimal — but consistent — probe rate; reported regret is therefore
an upper bound on what the sharpened schedule would achieve.
The empirical $d$-scaling slope $\approx-0.16$ in Figure~\ref{fig:d-scaling}
is closer to the sharpened $d^{1/3}=+0.33$ prediction than to the
original $d^{2/3}=+0.67$, providing post-hoc empirical support for
the variance-adaptive analysis.
\end{remark}

\begin{remark}[Combined LB/UB picture, after both extensions]
\label{rem:combined_after_ext}
Theorem~\ref{thm:spsc_lb_piecewise} and
Theorem~\ref{thm:spsc_regret_sharp} together yield
\[
\Omega\!\bigl(c^{1/3}\,K^{1/3}T^{2/3}\bigr)
\;\le\;\text{(probe-acquisition regret of SPSC)}\;\le\;
\tilO\!\bigl(d^{1/3}\,K^{1/3}T^{2/3}\bigr),
\]
tight in $K$ and $T$, with a residual $d^{1/3}$ gap that the
empirical $d$-scaling fit (Figure~\ref{fig:d-scaling}) suggests is
also loose; closing it likely requires the subspace-restricted
concentration sketched in Remark~\ref{rem:dim_dep}, which is beyond
the scope of this purely-additive sharpening.
\end{remark}

% =====================================================================
% End extensions.tex
% =====================================================================
